% SPDX-License-Identifier: Apache-2.0
%
% Vreteno: an RV32IMAC core in TxHDL, the first large design. Built by
% //docs:vreteno. One column, because it is mostly listings.
\documentclass[journal,onecolumn,9pt]{IEEEtran}

% One column: 80 columns fit at the body size, so nothing is reduced.
\newcommand{\listingsize}{\normalsize}
% SPDX-License-Identifier: Apache-2.0
%
% The house preamble, shared by every document in this directory. Loaded
% after \documentclass, before \title. Keeping it in one file is what
% stops two articles drifting apart in font, listing style or figure
% style.
% Computer Modern. IEEEtran selects Times (ptm) by default, so the face has
% to be chosen explicitly.
%
% Latin Modern is the Computer Modern variety used here, and the choice is
% forced rather than aesthetic. Under T1 encoding, plain `cmr` has no Type 1
% outlines unless cm-super is installed, and it is not in the pinned
% distribution, so pdflatex falls back to EC bitmaps and every glyph in the
% PDF becomes a Type 3 font. That was measured: the first build produced a
% document with 10 Type 3 fonts and no Type 1. Latin Modern is Computer
% Modern redrawn with a full T1 repertoire and real .pfb outlines, so the
% same design comes out scalable.
\usepackage[T1]{fontenc}
\usepackage[utf8]{inputenc}
\usepackage{lmodern}
% microtype is not in the pinned TeX distribution. emergencystretch alone
% keeps the two-column measure from producing overfull lines.
\setlength{\emergencystretch}{3em}

\usepackage{amsmath}
\usepackage{amssymb}
\usepackage{amsthm}
\usepackage{booktabs}
\usepackage{array}
% enumitem is not in the pinned TeX distribution; the list spacing below
% does what \setlist would have done.
\usepackage{float}
\usepackage{xcolor}
\usepackage{listings}
\usepackage{tikz}
\usetikzlibrary{arrows.meta,positioning,fit,backgrounds,calc}
\usepackage[hidelinks,bookmarks=true]{hyperref}

% Numbered, definition-style box for a rule the reader is meant to apply.
\theoremstyle{definition}
\newtheorem{rulebox}{Rule}
\newtheorem{example}{Example}

% Unnumbered asides.
\theoremstyle{remark}
\newtheorem*{tip}{Tip}
\newtheorem*{pitfall}{Pitfall}

% Tighter lists, in place of enumitem's \setlist.
\setlength{\topsep}{2pt}
\setlength{\itemsep}{1pt}
\setlength{\parsep}{0pt}

% Code in the running text. The typewriter face under T1 ligates `<<`
% and `>>` into guillemets, so a shift printed as \code{<<} came out as
% a French quotation mark (issue 571). microtype's \DisableLigatures is
% not in the pinned distribution, but the primitive it rests on is
% pdfTeX's own: \pdfnoligatures strips every ligature from the font it
% is given, and the current font inside \texttt is the typewriter face
% at the size in use, so each size is stripped the first time \code
% sets it. Code wants no ligature at all: two quotes are two quotes.
\newcommand{\code}[1]{\texttt{\ifdefined\pdfnoligatures\pdfnoligatures\font\fi #1}}
\newcommand{\term}[1]{\emph{#1}}

% Syntax highlighting. Four roles, distinguishable in colour and still
% distinguishable in greyscale, because an IEEE article gets printed: the
% keyword is the only bold role, the comment is the only italic one, and the
% two remaining roles differ in lightness as well as in hue.
\definecolor{kwblue}{RGB}{20,60,150}
\definecolor{cmtgreen}{RGB}{90,120,90}
\definecolor{strred}{RGB}{150,50,40}
\definecolor{numgray}{gray}{0.45}
\definecolor{framegray}{gray}{0.70}
\definecolor{bgshade}{gray}{0.97}
% A darker fill for figure cells. The listing background has to stay very
% light to keep code readable; a figure cell has to be clearly shaded or the
% distinction it draws is invisible in print.
\definecolor{cellshade}{gray}{0.90}

% The listing size is the document's choice, because it depends on the
% column: a two-column article sets `\listingsize` to 6pt and keeps its
% listings within 70 columns, or to 5.6pt when it includes source that
% rustfmt sets at 80, which is what fits a column's frame (issue 1061);
% a one-column document keeps the body size and includes source files
% at 80. `//docs:frame_test` checks every listing line against its frame. Lines are never broken; a line that does not fit
% overflows the frame, which is visible, rather than wrapping, which is
% not.
\providecommand{\listingsize}{\scriptsize}

% `'` is deliberately not a string delimiter: the language uses it for loop
% labels, as Rust does, so treating it as a quote would swallow the rest of
% the line.
\lstdefinestyle{txhdl}{
  basicstyle=\ttfamily\listingsize,
  keywordstyle=\bfseries\color{kwblue},
  commentstyle=\itshape\color{cmtgreen},
  stringstyle=\color{strred},
  numberstyle=\tiny\color{numgray},
  morekeywords={mod,use,pub,const,type,struct,enum,trait,impl,fn,async,let,mut,
    match,if,else,loop,while,for,in,break,continue,return,as,await,
    module,bus,channel,transaction,protocol,pipeline,flow,seq,config,
    port,stage,pipe,instance,connect,bind,tag,domain,blackbox,foreign,
    property,role,signal,handler,true,false,
    sequence,interface,modport,implements,package,import,var,wire,func,proc,
    case,default,next,fork,branch,join,state,fsm},
  morecomment=[l]{//},
  morestring=[b]",
  showstringspaces=false,
  columns=fullflexible,
  keepspaces=true,
  breaklines=false,
  % Framed, so a listing is bounded rather than floating in the column.
  frame=single,
  rulecolor=\color{framegray},
  framesep=3pt,
  backgroundcolor=\color{bgshade},
  xleftmargin=3pt,
  xrightmargin=3pt,
  aboveskip=6pt,
  belowskip=6pt,
  % Every listing is numbered and captioned, so it can be referred to by
  % number. `caption` without `float` numbers the listing and leaves it
  % where it was written, which is what a listing in running prose needs.
  captionpos=b,
  abovecaptionskip=2pt,
  belowcaptionskip=6pt,
}

% Figures drawn as text. Framed the same way, and with the highlighting
% switched off, because a box-and-arrow diagram is not code and half its
% words turning blue reads as an accident. Setting a style adds keys rather
% than clearing the ones already set, so the three roles are neutralised by
% hand rather than by leaving them out. Program output is printed in this
% style too, and a netlist's assignment can run past the frame, so a line
% that does not fit is broken and the rest indented; source never is.
\lstdefinestyle{ascii}{
  basicstyle=\ttfamily\listingsize,
  keywordstyle=\ttfamily,
  commentstyle=\ttfamily,
  stringstyle=\ttfamily,
  columns=fullflexible,
  keepspaces=true,
  breaklines=true,
  breakindent=2em,
  frame=single,
  rulecolor=\color{framegray},
  framesep=3pt,
  backgroundcolor=\color{bgshade},
  xleftmargin=3pt,
  xrightmargin=3pt,
  aboveskip=6pt,
  belowskip=6pt,
}
\lstset{style=txhdl}

% House TikZ styles. `shorten` lives on the arrow styles rather than on each
% arrow, so every arrow in the document gets the gap, including ones added
% later. TikZ clips a path at a node's border, so without it an arrowhead
% butts against the box with no gap at all. Keep `node distance` well above
% twice the shorten, or the arrow shrinks to nothing.
\tikzset{
  box/.style={draw,rounded corners=2pt,align=center,inner sep=4pt,
              font=\footnotesize},
  boxf/.style={box,fill=bgshade},
  lbl/.style={font=\scriptsize\itshape,align=center},
  fl/.style={-{Stealth[length=4pt]},shorten >=2.5pt,shorten <=2.5pt},
  fld/.style={fl,dashed},
}


% The contents list: the class gives a section's number 2.75em, which
% a Roman numeral past thirty overruns onto the title, so the box is
% wider here, and a subsection's entry starts where a title does.
\makeatletter
\def\l@section#1#2{\addpenalty{\@secpenalty}\addvspace{1.0em plus 1pt}%
    \@tempdima 5.75em \begingroup \parindent \z@ \rightskip \@pnumwidth%
    \parfillskip-\@pnumwidth {\bfseries\leavevmode #1}\hfil\hbox to\@pnumwidth{\hss #2}\par%
    \endgroup}
\def\l@subsection{\@dottedtocline{2}{5.75em}{5em}}
\def\l@subsubsection{\@dottedtocline{3}{10.75em}{5.5em}}
\makeatother

% SPDX-License-Identifier: Apache-2.0
% The boards' memory maps, written by //tools/memmap from the maps the
% designs decode (issue 444). A document asks for a table with
% \mmget{<what>}{<board>}, and a table the generator does not write
% stops the document rather than leaving a gap.
\input{tools/memmap/mm_tables}
\makeatletter
\newcommand{\mmget}[2]{%
  \@ifundefined{mm@#1@#2}%
    {\PackageError{memmap}{//tools/memmap writes no #1 for #2}%
      {Add it to tools/memmap/main.rs.}}%
    {\csname mm@#1@#2\endcsname}}
\makeatother


\InputIfFileExists{buildstamp.tex}{}{}
\providecommand{\buildstamp}{Draft build (outside the Bazel flow).}

% A range of the interrupt controller, by its begin/end markers.
\newcommand{\plicpart}[3]{%
\lstinputlisting[rangebeginprefix=//\ begin\{,rangebeginsuffix=\},%
rangeendprefix=//\ end\{,rangeendsuffix=\},includerangemarker=false,%
linerange=#1-#1,caption={#2},label={#3}]{lib/parts/src/plic.rs}}

% A range of the Sv32 unit, by its begin/end markers.
\newcommand{\mmupart}[3]{%
\lstinputlisting[rangebeginprefix=//\ begin\{,rangebeginsuffix=\},%
rangeendprefix=//\ end\{,rangeendsuffix=\},includerangemarker=false,%
linerange=#1-#1,caption={#2},label={#3}]{lib/parts/src/mmu.rs}}

\title{Vreteno: An RV32IMAC Core in TxHDL}

\author{Filip Filmar and Dragiša Janković%
\thanks{Both authors are independent. Filip Filmar is the
corresponding author, at f@hdlfactory.com; Dragiša Janković is
reached at linkedin.com/in/dragisa-jankovic.}%
\thanks{This document holds the Vreteno core, a RISC-V RV32IMAC processor
with machine, supervisor and user modes and Sv32 paging,
written in TxHDL under \code{//cpu/vreteno} of the project repository,
included verbatim at build time with the reference model it is checked
against, the programs it runs, and the trace and waveform the build
produced by running it. The language is described in the embedding
article~\cite{embedding}, the runtime in its own
document~\cite{runtime}, the smaller designs in the
examples~\cite{examples}, and the cover~\cite{cover} lists all of them.}%
\thanks{The authors used a large language model, Claude, as an
assistant in exploring the concepts and in writing and constructing
the documents and the programs; every commit of the repository records
this attribution and includes the exact prompts that guided it.}%
\thanks{\buildstamp}}

\markboth{Vreteno: An RV32IMAC Core in TxHDL}{Filmar and Janković: Vreteno}

\begin{document}
\maketitle

\begin{abstract}
Vreteno is a 32-bit RISC-V core, the RV32I base integer set, the M
extension, the A extension's atomic instructions and the C extension's
compressed instructions, with machine, supervisor and user modes and
Sv32 paging, written in
TxHDL: the first design large enough to ask whether the language
scales past its examples. This is its fifth form: a three-stage
pipeline, one process in which a fetch stage reads the instruction
memory into an instruction register, a compressed instruction
expanded there into the one it stands for, an execute stage runs the
instruction before it and reads the data memory into a register at
its edge, and a writeback stage extends a load, writes the register
file and retires, with the one hazard closed by forwarding, except
after a load, where it is closed by a stall of one cycle, and one
bubble per taken branch or jump, written in the subset the lowering
reads, so that the same file simulates and is a netlist, and with
machine-mode traps: an environment call or a word the core does not
know saves its address and its cause in the control registers and
enters the handler, which returns through \code{mret}; and with
the M extension on one sequencer, a multiply one step in the part's
multipliers and a division thirty-two, which stalls the instruction
in execute; with an interrupt line, taken in place of the
instruction in execute when \code{mie}, \code{mip} and
\code{mstatus} allow; and with a bus, a request channel out and a
response channel back through a router, on which the data memory, a
timer and a serial port are units of their own, the timer the second
interrupt, and the port where the demonstration says OK and echoes
what a terminal types back. It is checked twice, and it has run on
the board. In lockstep against a reference model of the instruction set
written as a plain program: the model steps when the core retires an
instruction, and after every cycle the program counter, the
thirty-one registers, the control registers and the halt must agree,
and at the halt the data memory must; the demonstration program and
sixty-four random ones pass, and a deliberately broken shift, a
branch that does not squash the word after it, and a return from a
trap that leaves the wrong status, each fail at their first use. And as a netlist: the VHDL
and the Verilog the lowering writes, with the program in its
instruction memory, are simulated under nvc and Verilator against the
trace the Rust core wrote, every register and output at every tick,
and agree. Every listing here
is the file in the tree, and the trace and the figure are what the
build produced when it ran the core.
\end{abstract}

\begin{IEEEkeywords}
RISC-V, RV32IMAC, processor, hardware description languages, Rust,
lockstep testing.
\end{IEEEkeywords}

{\footnotesize\tableofcontents}

\section{What Vreteno Is}
\label{sec:v-what}

\emph{Vreteno} is the Serbian word for a spindle: the part that turns
and on which the rest is wound. The core is the first design in TxHDL
with more than a handful of registers, and it exists to find out three
things: whether the language's constructs, a unit with registers and
memories, one process, waits, reads and deferred drives, are enough to
write a processor without stepping outside them; whether a design of
this size can be tested to a standard the small examples cannot reach;
and what the lowering must grow to before such a design becomes a
netlist.

The plan had three steps, and this document covers all three and a
fourth. The core was written first as a single-cycle machine in
simulation form and checked against a reference model, so the design
was known correct before anything else was asked of it. Then it
became a pipeline of two stages, and the lockstep check validated
that refactoring: the test changed in one place, the core in
one file, and the same programs passed. Last, the lowering
grew until the core was a netlist, and the core was rewritten in the
subset the lowering reads, with the lockstep check again holding the
rewrite to the same behaviour; and the netlist simulates, under nvc
and Verilator, against the trace the Rust core wrote, which is the check the smaller
examples~\cite{examples} pass. What the lowering had to grow for
this is stated in Section~\ref{sec:v-netlist}. Then, with the
netlist synthesized and its memories and its timing known, the core
became the pipeline of three stages it is now, so that the data
memory could be a block RAM, and the lockstep check held that
refactoring too.

\section{The Instruction Set as Bits}
\label{sec:v-isa}

Listing~\ref{lst:v-isa} is RV32IMAC as encoders and a decoder. An
encoder per mnemonic builds the word, the six formats being six
functions of their fields; the decoder takes a word apart into its
mnemonic, its registers and its immediate, sign-extended. The core
does not use this decoder: it decodes on its own, in
Section~\ref{sec:v-core}, from the fields of the word with the
language's \code{slice}, \code{concat} and \code{sext}. So there are
two decoders, written differently, and the lockstep test compares
what they lead to; a defect in either shows up as a disagreement. The
disassembler is for traces a person reads, such as
Listing~\ref{lst:v-out}. The system instructions are there with the
rest. \code{ecall}, \code{ebreak} and \code{mret} share an opcode
and a function code and differ in their upper twelve bits; in the
six CSR forms those bits name the register. The eight machine-mode
registers a trap uses, \code{mstatus}, \code{mtvec},
\code{mscratch}, \code{mepc}, \code{mcause}, \code{mie},
\code{mip} and \code{mtval}, and the causes the core raises, are
constants beside them, with the external and the timer interrupt's
bits and causes, and the timer's and the serial port's addresses.
The two machine counters are there, \code{mcycle} and
\code{minstret}, each 64 bits and each read as two words. They are
what a profiler wants, and what lets the core's own numbers be
measured by a program running on it rather than only from the
simulation around it. \code{mcycle} counts while the core runs and
stops when the halt stops it, since nothing restarts this machine and
a count that ran on afterwards would be a number nobody could read.

They are the one part of the core the lockstep test does not check,
by architectural design. The model steps when the core
retires, so it knows how many instructions have gone by and cannot
know how many cycles the pipeline spent; a program that read
\code{mcycle} into a register would get a number no model could
produce without knowing the microarchitecture, which is the thing the
model exists not to know. So the counters are checked by a run of
their own, \code{//cpu/vreteno:counters\_test}, where a program reads
both around a division and finds the cycles many more than the
instructions.

\code{wfi} is there as well, and it waits rather than spinning. The
specification allows either, and a core that returned at once would be
correct; this one stops fetching until an interrupt is pending and
enabled, and the cycles between are cycles in which it retires nothing.
The test measures that silence rather than trusting it: a spin would
retire something in nearly every cycle.

What the wait does not ask is whether interrupts are enabled globally.
The specification lets a core resume with \code{mstatus} closed, and
that is the case that matters, because a kernel idles inside its own
interrupt lock: a wait that waited for the global bit would never end.
So the line coming up ends the wait, and whether the trap is then
taken is the usual question, answered the usual way.

\code{mhalt} at \code{0x7c0} is this core's own, in the space the
specification leaves to a design: it reads as zero, and a write of an
odd value to it stops the core when that instruction retires, which
is what \code{halt} encodes and what a program ends with.
\code{dcsr} at \code{0x7b0} and \code{dpc} at \code{0x7b1} are the
debug specification's two, and Section~\ref{sec:v-core} says what
the core keeps in them. \code{mbusquiet} at \code{0x7c1} is the
core's own, beside \code{mhalt}: its low bit set, a load the bus
refused reads the zero the bus answered and a store it refused is
dropped; clear, which reset leaves, either is an access fault, as
Section~\ref{sec:v-core} says.

Five more say what the machine is rather than what it is doing.
\code{misa} reports RV32IMAC with supervisor and user mode, and \code{mvendorid}, \code{marchid},
\code{mimpid} and \code{mhartid} read as zero: an unassigned vendor,
an unspecified architecture and implementation, and one hart, which is
hart zero. They are there because a stock kernel reads one of them
before it does anything else. Zephyr's reset code begins
\code{csrr a0, mhartid}, and until the core answered it that
instruction was an illegal one, so the kernel's first act was a trap
into a handler it had not installed.

The four that begin \code{0xf} are read only, and the specification
means that strictly: a write to one is an illegal instruction rather
than a write that is ignored. The core raises it, and only for an
instruction that really writes. \code{csrrw} and \code{csrrwi} always
do; a set or a clear writes only when its source field is not
\code{x0} or a zero immediate, which is stated in terms of the field
and not of the value the register happens to hold. A read is a set
with \code{x0}, so the kernel's first instruction is answered.
\code{misa} is not among them: the specification makes it writable and
lets an implementation ignore what is written, which is what this one
does. The M extension is the
register format with \code{funct7} one: the four multiplies, which
differ in which operand is signed and which half of the product they
keep, and the two divides and two remainders, signed and unsigned. The
A extension's word forms are opcode \code{0x2f} with \code{funct3}
two, told apart by the five bits at the top: \code{lr.w}, whose
\code{rs2} must be zero, \code{sc.w} and the nine AMOs, swap, add,
the three logical ones, and the signed and unsigned minimum and
maximum (issue~1010). Bits 26 and 25, \code{aq} and \code{rl}, are
decoded and set aside, since one hart in order already orders them.
User and supervisor mode (issue~1012) add \code{sret}, a system
instruction beside \code{mret}; the supervisor's registers,
\code{sstatus}, \code{sie} and \code{sip}, which are views of the
machine's, and \code{stvec}, \code{sscratch}, \code{sepc},
\code{scause}, \code{stval} and \code{satp}; the delegation registers
\code{medeleg} and \code{mideleg}; the counter enables
\code{mcounteren} and \code{scounteren}; and the unprivileged
counters \code{cycle}, \code{time} and \code{instret} with their high
halves. A CSR's address names in its bits 9 and 8 the least privilege
that may reach it, and one whose top two bits are set is read-only.
A unit test checks each format's encoder against the decoder.

The compressed instructions are sixteen-bit spellings of instructions
above, and \code{isa} has an encoder for each, named \code{c\_} and
the assembler's mnemonic. \code{compressed} reads a halfword as the
thirty-two bit instruction it stands for, or refuses it: all zeros,
\code{c.jr} of \code{x0}, \code{c.lwsp} into \code{x0}, \code{c.lui}
and \code{c.addi16sp} of zero, the floating-point and 64-bit forms,
and shifts by 32 or more, which RV32C leaves to custom extensions. A
hint, \code{c.addi} of zero or \code{c.mv} into \code{x0}, is an
instruction that does nothing, and expands to the one it is spelt
as. \code{compress} goes the other way, as an assembler does when it
may, and a test checks that every instruction a halfword stands for
compresses to a halfword that expands back to it.

This decoder has a checker from outside the project. The build
assembles all 49\,152 halfwords whose low two bits are not both set
with the GNU assembler of the toolchain it already fetches and lists
them with its disassembler, and \code{//cpu/vreteno:compressed\_test}
spells each instruction the listing names with the thirty-two bit
encoders and requires \code{compressed} to give the same word, and
to refuse every halfword the listing does not name. The two differ in
two places, where the test follows the specification: the
disassembler names \code{c.addi16sp} of zero and shifts by 32 or more.
Of the 49\,152, 28\,823 are instructions.

\lstinputlisting[caption={\code{isa.rs}: the encoders, the decoder,
the disassembler.},label={lst:v-isa}]{cpu/vreteno/src/isa.rs}

\section{The Reference Model}
\label{sec:v-model}

Listing~\ref{lst:v-model} is the instruction set as a program. Its
state is the architectural state and nothing else: the program
counter, the registers, a data memory, the eight control registers,
and whether it has halted. One \code{step} is one instruction,
written against the decoder in the plainest Rust that says what the
manual says. It is the oracle: the core is right when it agrees with
this, cycle by cycle. The data memory begins at \code{DATA\_BASE}.
Below it a load reads the boot memory and above it the devices
answer; a load or store in the gap between the two memories is a
fault, so a wild address in a random program is caught rather than
wrapped.

The traps are the manual's, machine mode only. An \code{ecall}, or a
word the decoder does not know, is a trap: the address of the word
goes to \code{mepc}, its cause to \code{mcause}, eleven for the
call and two for the illegal word, the interrupt enable of
\code{mstatus} is saved into its previous-enable bit and cleared,
and the next instruction is at \code{mtvec}. \code{mret} is the
way back: the counter takes \code{mepc}, the enable takes the saved
bit, and the saved bit is set. A CSR instruction reads the register
its upper twelve bits name into \code{rd} and writes it, replaced,
set or cleared by a register or by a five-bit immediate; a name the
model does not have is an illegal word, and traps. \code{mtval}
takes the illegal word itself, the address an unaligned access asked
for, the breakpoint's own address, and zero for any other trap. The
external interrupt is a bit in \code{mip}, which \code{line} sets
to the line at each edge, as the privileged specification has it: the
interrupt controller sets and clears it, and a write to \code{mip}
changes nothing; when it is set,
enabled in \code{mie} and interrupts are enabled in
\code{mstatus}, the next \code{step} traps instead of executing,
with the interrupt's cause, and returns to the same address. The
controller is five words in the window at \code{CLINT\_BASE}, the
count and the compare, each in two halves, and \code{msip}, read and
written as memory is; the timer's pending bit and the software one
are its two lines, read in \code{mip} but not stored there, and the
order they are taken in is the specification's: the external
interrupt, then the software one, then the timer's. The model has no
clock and no bus, so what a device answers is not its own. The caller
sets \code{dev\_word} before a step to what the bus answered the
core's load from a device, and \code{tirq} to the line as the core
saw it; the data memory's words the model keeps itself, and the
memory unit's lanes are compared with them at the halt. A store to a
device is the device's business; the model keeps what it can check
at the end, the timer's compare and the bytes given to the serial
port. Which
interrupt that step takes, if any, is the caller's word too, since
the core decides on the bits as they stood a cycle earlier, and the
lockstep harness propagates the core's decision across. Only the write to
\code{mhalt} halts. The fetch reads a halfword, the upper half of its
word when bit 1 of the address is set, and takes a second when the
first's low two bits are both set; a compressed instruction is read
as the one it stands for, the counter moves on by the instruction's
length, and a jump links that address. A halfword that is no
instruction decodes as illegal, and is the trap value. The
M extension is Rust's arithmetic on the
right types: a product's low word is the same whatever the signs, the
high word is the top of a 64-bit product of the operands extended
as the instruction says, and a division is the wrapping one, which
gives the manual's answer for the one overflow; division by zero,
which Rust would refuse, is spelled out, all ones and the dividend.
The A extension is a read, the operation and a write, with the old
word in \code{rd}; \code{lr.w} reserves the word and \code{sc.w}
stores only while it is reserved, writes 0 if it stored and 1 if not,
and uses the reservation up either way. A misaligned one traps,
\code{lr.w} as a load and the rest as a store, and a word the bus
refuses is a load's access fault for \code{lr.w} and a store's for
an AMO. The privilege is a number, 3 machine, 1 supervisor, 0 user. A
trap below machine mode whose cause machine mode delegated, in
\code{medeleg} for an exception and \code{mideleg} for an interrupt,
goes to the supervisor through its registers, saving \code{SIE} and
the mode in \code{SPIE} and \code{SPP}; any other goes to machine mode,
saving \code{MIE} and the mode in \code{MPIE} and \code{MPP}. An
interrupt machine mode keeps is taken below machine mode always and in
it with \code{MIE}; one it delegated is taken below supervisor mode
always and in it with \code{SIE}; the machine's come first. The
model takes the supervisor's external line on \code{sline}, beside
\code{line} for the machine's.
\code{mret} and \code{sret} return to the saved mode and leave user
mode in its place, and \code{ecall}'s cause is 8 plus the mode it came
from. A CSR the mode may not reach, \code{mret} below machine mode, and
\code{sret} and \code{wfi} in user mode are illegal instructions.

The same model is also a fast machine, for software too large for the
cycle-level runs: a kernel, and what boots it (issue~\#1016). Given a
\code{Bus}, it sends every fetch, load and store there instead of to
its own memories, and a store goes out as a word and a mask of its
lanes, so that a device is never read on the way to being written.
\code{cpu/vreteno/src/machine} puts the board around it at the
addresses its maps give: the DDR3 as a gigabyte of plain memory, the
CLINT with its count moving one an instruction, the PLIC with its two
targets, and the serial port as SiFive's. A loader lays an image and
a device tree blob into the DDR3 and starts the hart as a bootloader
would, \code{a0} the hart's number and \code{a1} the blob's address,
and \code{//cpu/vreteno:machine} runs it with the serial port's output
on the terminal and the terminal's input on the port, so a shell on the
machine can be typed into or fed from a pipe (issue~\#1127). A program built for the board's DDR3 runs there
unchanged: \code{hello\_ram\_bin} prints its line and halts after 267
instructions. Without a bus the model is the board's reference as
above, and the lockstep test holds the core to it as before.

What Linux boots from is one image, which fastboot loads at
\code{0x4000\_0000} and jumps to (issue~\#1019). \code{//tools/bootimg}
lays it out: OpenSBI's \code{fw\_jump} at \code{0x4008\_0000}, the
device tree at \code{0x4030\_0000}, the kernel at \code{0x4040\_0000},
where OpenSBI jumps, and the initramfs past the kernel's \code{\_end},
taken from its \code{System.map}, so that the kernel's BSS is clear of
it. The image does not hold them there, though. Placed at those
addresses it was half zeros, and fastboot sends every byte at about
26\,KB/s, so the parts are packed one after another and the initramfs
gzipped, which the kernel unpacks itself (issue~\#1201). A few
instructions at the front stand in for a bootloader. First they move
each part to its address, the last part first and each a word at a
time from its top down: no part is packed above its address, so no
word is overwritten before it is copied. Then \code{a0} is zero,
\code{a1} the tree, and a jump to OpenSBI. Before the jump
they leave the serial port and the interrupt controller as SiFive's
parts leave them at reset, as Linux's drivers expect
(issue~\#1136). Our port resets with its receive interrupt enabled,
so the serial loader's bytes raise its line, and the controller keeps
the request. Linux's driver enables the source before it has
registered the port, so the request reached the handler with no port
and the kernel panicked. The shim writes the port's \code{ie} to zero,
claims and completes every request the controller holds, and leaves
every source disabled at priority zero. \code{boot\_loaded\_test}
boots Linux on the machine with the port as the loader leaves it. The
tree is
written last, with the initramfs's range and the command line in its
\code{/chosen}, and the packer refuses an image longer than fastboot
takes, its download size less the header page the host puts in front
of a file, naming the sizes. The command line is
\code{earlycon console=ttySIF0} unless a boot asks for another, so the
kernel speaks on the serial port from its first line (issue~\#1125).
The rule \code{boot\_image} chains the steps, and
\code{fastboot boot NAME.bin} sends the result; its \code{model}
attribute adds \code{mem=64M} for a boot on the machine below, where
setting up the board's whole gigabyte is most of the kernel's early
work, and the board's own image boots with all of it. The
machine above is its first reader: a test boots a packed image there
and checks the shim's registers and every part's address before an
image ever reaches the board.

Linux reaches the network through its own LiteEth driver,
\code{litex\_liteeth}, on the Ethernet port (issue~\#1203). The
port's registers are LiteEth's at LiteX's offsets, and its two
receive and two transmit slots are the DDR3 buffers the tree reserves,
so the driver runs as the kernel ships it. The tree names the port as
\code{litex,liteeth}, with its registers and its buffers. The machine
has the port too: its registers by the hardware's own map, its engines
as copies to and from the DDR3, and with \code{--eth-peer} a station
on the cable that answers ARP and ping at 10.0.0.2. It has the entropy source too, its registers by the part's map and
a fixed sequence in place of the rings, so that a Zephyr image whose
network stack asks for entropy boots there: \code{fastboot\_model\_test}
boots the fastboot server to its listening line (issue~1386). With \code{--fastboot-peer} the far end of the cable is a fastboot
client instead, smoltcp's TCP/IP with the protocol on top, and it
sends the server a download, which arrives through the port's slots
at the line's pace and is dropped as the board drops it (issue~1390).
\code{fastboot\_download\_model\_test} sends 256~KiB, checks the
staged bytes against the image, and says the steps a byte took, the
retransmits and the dropped frames. A step is an instruction, since
the machine has no memory latency, so the number is the work the
receive path does and how TCP fares on the port, and the board still
says what the DDR3 costs. Its frames arrive
at the line's pace and wait for the driver as the hardware's do, at
most one a slot, and a frame that finds both slots held is dropped and
counted in \code{rx\_errors} (issues~1313 and~1314), so a driver
too slow for a burst loses frames on the machine as it does on the
board.
A step of the machine is an instruction, which is what lockstep
compares, and \code{mcycle} reads zero there. With \code{--timing} a
step is charged what the core would spend on it instead (issue~1392):
one cycle an instruction and four more when control does not fall
through, and for a data access forty more from the DDR3, three into
it, nine from a memory on the bus and one from the core's own, eight
from a peripheral and two into one. \code{mcycle} and the timer count
those cycles. The costs are fitted to the board's own programs, and
\code{timing\_test} holds the model to the fit: \code{cpi.rs}'s warm
loop takes 7000 cycles where the board took 7036 from the data memory
and 7088 from the DDR3, and \code{ethperf.rs} copies a full frame into
a DDR3 slot in 20\,142 where the board took 19\,926. A cold run misses
by the instruction cache's fills, which the model does not have, and a
transmit by the engine, which it does not time.
\code{boot\_net\_test} boots an image whose command line asks
\code{/init} for the network, and passes when Linux has pinged that
station and had its replies.
\code{/init} reads two words of the command line for it (issue~1318):
\code{txhdl.ip=} brings \code{eth0} up at an address, and
\code{txhdl.ping=} pings a host and says whether it answered;
\code{txhdl.net} is the two at the machine's addresses.
The board's image specifies \code{txhdl.ip=192.168.1.50/24}, the
address the fastboot server has, so on the board Linux is on the
network beside srv's 192.168.1.1 with nothing typed.
\code{boot\_ip\_test} boots the machine's image with the two words
spelled out.

\lstinputlisting[caption={\code{model.rs}: the reference
model.},label={lst:v-model}]{cpu/vreteno/src/model.rs}

\section{The Core}
\label{sec:v-core}

Listing~\ref{lst:v-core} is the core. It is one unit whose state is
the fetch stage's program counter; an instruction register with the
program counter it was fetched at, a valid bit, and a bit that says
the instruction was compressed, which together
are the boundary between fetch and execute; the writeback registers,
the boundary between execute and writeback, which hold whether an
instruction is there, its program counter, its destination, its
value if it is not a load, the raw word read for a load with the
lane and the width to take from it, the word itself, and whether it
halts; the halt itself; the eight control registers; the sequencer's
seven, a busy bit, a count, the two halves of the pair it shifts, the
multiplicand or divisor, and two sign bits; the word a load read
from the bus and whether a load is waiting on it; and two memories:
the register file, and the instruction memory loaded from the
program with \code{Mem::with}. The data memory is a device on the
bus, a unit of its own. Its inputs are
the reset, the three interrupt lines, the bus's answers and the
debugger's requests, and its outputs are the halt, the word
executed, a \code{Writeback} value, the bus's requests and the
debugger's answers; the \code{Writeback} value says whether an instruction retired this cycle,
the register it wrote and what it wrote, which is what the waveform
decodes field by field and the lockstep test steps the model on.
The pieces of the step that are functions of their operands alone
come first in the file, each a function under \code{\#[lower]} with
a name and a sentence of its own: the expander of compressed
instructions and the encoders of the formats it builds, the five
immediates, the ALU, the
branch condition, a loaded word's extension, what a store puts on
each lane and which lanes it writes, the CSR read, the CSR number
check and the CSR's new value, the sequencer's signs and its result.
The step calls them, and the lowering writes them in where it does,
so the step is the pipeline's decisions and drives, some four
hundred and seventy lines, and the netlist is the one it made when
the pieces were written in place.

The process is one \code{loop} with one wait, the rising edge, and
everything after the wait is one cycle of all three stages at once.
The fetch stage reads the instruction memory at its program counter
and at the word after, since an instruction may cross from one word
into the next. The halfword at the counter is the upper half of its
word when bit 1 of the address is set. If its low two bits are not
both set it is a compressed instruction, and \code{expand} makes it
the thirty-two bit instruction it stands for; otherwise the halfword
after it is the instruction's upper half. The counter moves on by two
or by four, and the instruction register takes the instruction with
a bit that says which, so nothing after the fetch knows there are two
lengths but the link a jump writes, two bytes on for a compressed
call. \code{expand} is one \code{select!} on the three bits of the
major opcode above the two that say compressed, over the instructions
the encoders build from the halfword's fields, and it gives a halfword
that is no instruction back as itself: no thirty-two bit instruction
has its low two bits clear, so it traps as illegal, with itself as
the trap value, which is what the specification asks. A test compares
it with \code{isa::compressed} on every compressed halfword. The
execute stage takes the instruction register:
the decode, the fields of the word by \code{slice} and the five
immediates assembled by \code{concat} and \code{sext}, as the manual
draws them; two plain reads of the register file; the execute, an
ALU shared by the register and immediate forms and a branch
condition, each a function of its inputs; and the store and the
load, which go out on the bus as one request, the address, the data
in its lanes and the lanes a store covers, a load's word coming back
on the bus into the writeback stage's register. The writeback stage
takes the byte or the half out of that word by the lane and the
width the load specified, extends it, and writes the register file, or
writes the value execute computed; an instruction retires there.

The one hazard a third stage brings is closed in one place. An
instruction in execute may read a register the one in writeback has
not yet written; the operand reads check the writeback stage's
destination and take its value instead, the forwarding path, and
since the register file is written at the end of writeback, an
instruction two behind reads what it wrote. Register zero reads as
zero before that check, so a writeback to nowhere forwards nothing.
A load is the exception. Its word lands in the writeback register at
the edge and is extended there. That value, forwarded through the
ALU and the branch compare into the fetch counter, is the longest
path in the core, and once the traps had lengthened the fetch's
multiplexer it no longer fit in ten nanoseconds. So an instruction
that reads what the load in writeback will write waits one cycle,
the fetch holding with it, and reads the register file, which has
the value by then. The stall is one \code{let} and the one case in
which the core waits. Only the execute stage's values are forwarded,
and those are registers, and fast.

The traps are in execute, beside the ALU. The system opcode is
taken apart by its function code and its upper twelve bits into the
CSR operation, the call, the break and the return, and the
\code{select!} that says whether the core knows the opcode has an
arm for it, so an unknown system word is an unknown word like any
other. A trap is a call or an unknown word in a valid instruction;
its drives write the address, the cause and the new status into the
control registers, its target is \code{mtvec}, and it is a redirect,
so the word behind it is squashed as a taken branch's is, and it
writes no register. \code{mret} is a redirect to \code{mepc} with a
drive of the status. A CSR instruction selects the old value by the
register's number, computes the new one by the function code, the
source a register or the immediate in the \code{rs1} field, and
drives the one register named, with the old value as its result
through the ordinary writeback. The model's CSR step is one function
over an enum; the core's is five entries of one \code{with!} under
one enable, which is what five registers with one write each are.

The M extension is a sequencer, and the second use of the stall. A
multiply or a divide in execute starts it: the operands' magnitudes
go into the pair and the multiplicand or divisor register, the signs
into two bits, and the instruction stalls, the fetch holding with it.
A multiply is one step: the 64-bit product of the two registers,
\code{mul::<64>} of the examples document's multiply-accumulate
\cite{examples}, into the pair, which synthesis makes of the part's
DSP blocks. A division is thirty-two steps, one a cycle, each shifting
the pair left, subtracting the divisor from the high half when it
fits, and shifting the fit in as the quotient bit, which leaves the
quotient in the low half and the remainder in the high. When the
count is up the instruction runs: the result is the pair with the
signs put back, the product or quotient negated when the signs
differed, the remainder given the dividend's sign, a quotient by
zero left as the all-ones the steps produce, and it goes through the
ordinary writeback. So a multiply takes three cycles and a division
thirty-four, and one set of registers serves both. The multiplier
reads registers and writes registers, so it is a path of its own, a
cycle long, and not on the one that sets the clock; the first form
of the sequencer multiplied by thirty-two shifts and adds, and the
change to the multiplier was one step's arm and the count it starts
from.

The A extension (issue~1010) is the third use of the writeback stage's
hold, and it is placed to keep its arithmetic off the two paths that
set the clock, the operands' forwarding into the adder and the branch
into the next counter, which the design note \code{docs/sv32-timing.md}
measured. An AMO goes out as a load does, from execute, and its word
lands in writeback. There the instruction stays for two more cycles:
in the first the new word is computed by \code{amo\_alu}, from two
registers, the old word and the register operand kept from execute;
in the second it goes out as a store, from writeback, with the
address the load used. Then the instruction retires with the old
word. \code{lr.w} is a load that also sets a one-word reservation,
and \code{sc.w} is a store that goes out only while the reservation
holds, writing 0 or 1. The reservation is compared with the address
register itself rather than with the adder's sum, which an A
instruction does not need since it adds no offset. Every A
instruction first waits, as \code{fence} does, for the stores already
posted, so \code{aq} and \code{rl} need nothing more. The cost is
time, not clock: an AMO on the data memory takes about forty cycles,
most of them the wait for the previous one's store.

User and supervisor mode (issue~1012) are a two-bit privilege register
and the supervisor's registers beside the machine's, and they are
built so that a trap's destination is decided the way its cause is:
\code{to\_s} is one bit of \code{medeleg} or \code{mideleg}, picked by
the cause's low five bits, and it chooses the vector, \code{stvec} or
\code{mtvec}, the registers the trap writes, and the status it
leaves. The interrupt is chosen from the pending and enabled bits
split by \code{mideleg}, the machine's set first when the mode allows
it, which is a function of registers alone and so adds nothing to
the paths that set the clock. \code{mip}'s supervisor bits are
software's: machine mode sets them, which is how a supervisor gets
its timer from the SBI, and the supervisor clears its own software
interrupt through \code{sip}. \code{time} reads the timer's count,
which the timer hands the core on a port of its own, so that the CLINT
and the counter cannot disagree. Writing the lockstep test for this
found that a \code{wfi} the core refused in user mode still set its
wait, and a core that waited for an interrupt nothing would send; the
wait now asks that the instruction be legal.

Three of the cycle's decisions, whether the instruction in execute
is stalled, whether an interrupt takes its place and whether the
fetch is redirected, are wires kept as fields, \code{Wire<Bit>} of
the runtime: the process sets each and reads it back in the step, the
trace shows them under the core's name beside the registers, and the
netlist declares them from the fields, so the testbench that checks
the core checks them too. Any \code{let} of the loop can be made
one, at no cost in hardware, which is how an internal signal is
looked at.

The interrupt is the second input of the core, and it is taken in
execute as a trap is. The pending bit of \code{mip} is the line,
taken at each edge, and a write to \code{mip} changes nothing, since
the specification makes the bit read only and leaves it to the
interrupt controller. It once latched until software cleared it, which
software written to the specification never does: Zephyr took one
interrupt, returned to the bit still set, found nothing to claim, and
halted (issue 788). An instruction in execute that is not
stalled is live; when the pending bit, its enable in \code{mie} and
the enable in \code{mstatus} are all set, the interrupt takes the
live instruction's place: it runs nothing, its address goes to
\code{mepc} with the interrupt's cause, \code{mtval} takes zero,
the status is saved and cleared as for any trap, and the fetch is
redirected to \code{mtvec}. The decision is made a cycle before it is
acted on (issue 1331): whether one is pending and enabled, for which
mode, and from which set, go into registers, and the take reads them,
since the enables' logic in front of the fetch's address was the
flagship's worst path. A cycle later is nothing a program can see.
What it could see is an interrupt taken just after it was forbidden,
so the decision is dropped after a write of \code{mstatus},
\code{sstatus}, \code{mie}, \code{sie}, \code{mideleg}, \code{mip} or
\code{sip}, an \code{mret} or \code{sret}, a trap, and debug mode.
Only as the instruction runs: a CSR access merely sitting in execute
once dropped it every cycle, and a loop reading \code{mip} while it
waited for an interrupt never took one. The instruction still passes through
the writeback stage, writing nothing, so that it retires as a trap
does and the lockstep test steps the model on it. A division in
flight when an interrupt is taken is cancelled, and starts again
when the handler returns, on the registers as they are then: the
sequencer must never step under an instruction other than its own,
which is a defect the random programs found. On the board the line is the
interrupt controller's, whose input from the board itself is tied
low; the run and the test drive it.

Debug mode is the third thing that keeps a live instruction from
running, and it is entered as an interrupt is taken: before the
instruction, whose address goes to \code{dpc} and which runs when
the mode is left. Three things ask for it, and the cause field of
\code{dcsr} says which did: a halt request on the core's
\code{haltreq} input, the instruction after a single step, and
\code{ebreak} when \code{dcsr.ebreakm} is set, which otherwise
raises the breakpoint trap. In the mode the fetch is held, nothing
runs and nothing retires, and an interrupt that arrives waits with
its pending bit set. A resume, \code{resumereq} held high for a
cycle and acted on once per entry, redirects the fetch to
\code{dpc}; with \code{dcsr.step} set it arms a step, and the next
instruction to run asks to enter again before the one after it,
which is how a debugger walks a program an instruction at a time.
A core that has stopped itself, by writing \code{mhalt}, enters the
mode on a halt request too, with \code{dpc} the instruction after the
write, and the stop ends, so a debugger can read what a finished
program left and resume it (issue~1047).
\code{dcsr} holds the specification's version, the privilege,
\code{ebreakm}, \code{step} and the cause, and nothing else is
written. The two requests are inputs of the core and the mode is
its output \code{dbg}; on the board the debug module of issue 154
drives them, and the lockstep test drives them from a plan,
Section~\ref{sec:v-lockstep}. The CSR instructions name \code{dcsr}
and \code{dpc} only in debug mode, as the specification has it: from
machine mode a read or a write of either is an illegal instruction
and writes nothing (issue~972). A debugger sets \code{step} and
\code{ebreakm} through the module's own port, the abstract command,
while the core is halted.

The reset line puts the CSRs back as configuration left them: the
status, the trap vector, the enables and the pending bits, the trap
registers, the scratch register and both counters, and debug mode
with its two registers. The core declares \code{rst} itself, so no
register of its is put back by the netlist, and until the line did
this a program started by it ran with the previous program's trap
vector and interrupt enables (issue 419). The specification asks only
that interrupts be disabled at reset; the rest is this design's
decision, taken so that a program started by the line sees what a
program started by configuration sees. The register file is not
reset, which is what a register file is, and the lockstep test relies
on that: its program counts its own starts in a register.

The bus is everything from the data memory up in the address map,
and it is two channels: a request out of the core, one word of the
address, the data in its lanes, the lanes a store covers and whether
it is a store; and a response back, the word read. A channel is
what the language has for two units, registered on both sides and
indifferent to which ran first in the step, so a device is a unit of
its own, and the core's process talks to it by predicate, as the tap
of the examples document does~\cite{examples}: a receive with no
wait before it, which takes the response whenever it comes, and a
send under \code{if}, whose \code{valid} is the condition. A
store to the bus goes out when the bus has room and the instruction
runs on; a load goes out, sets the wait bit and moves on to
writeback, which holds it, and the whole pipeline behind it, until
the response has landed in a register of its own there. The hold is
the third use of the stall,
and it is a register's say: an earlier form stalled the load in
execute on the decode of its address, and that decode, the address
add and the fetch's multiplexer were the path that limited the
clock. A load or a store also waits for room on the bus whatever its
address, so that the fetch's hold never hangs on the address; the
devices keep up, so there is always room. A load that has gone out
is not taken by an interrupt: the interrupt waits for the next
instruction, so a device is never read twice for one load. A store
is posted, and a program that needs it to have landed before it
goes on says \code{fence}: the core counts the stores it has posted
and not yet had answered, and a fence waits in execute until the
count is zero. A load already waits for its answer, so that is the
whole of what a fence has to wait for; until issue 432 the
instruction retired at once and ordered nothing.

The bus says how an access went, and the core listens. A load the
bus refused, by a peripheral that failed or by the router for an
address that is nobody's, retires as a trap: the load access fault,
cause 5, with the address in \code{mtval} and the load's own address
in \code{mepc}, nothing written, the instruction behind it squashed
and the fetch restarted at the handler. A store is posted, so its
refusal comes back later; the core keeps it and raises the store
access fault, cause 7, before the next instruction to run, whether or
not interrupts are enabled, with \code{mepc} naming that instruction
and \code{mtval} zero, since the address is gone. Until this a load
from nowhere read zero and a program proceeded as though it had read
a word, which two board programs paid for in silence (issue 417).
The remote peripheral refuses on purpose when its patience runs out,
so a program that means to poll it sets \code{mbusquiet} and reads
the zeros as before. A fetch the bus refused comes back as a zero
word with a mark, and the instruction taken from it raises the
instruction access fault, cause 1, with its own address in
\code{mtval}, rather than the illegal instruction the zero would be
(issue 423).

The bus is AXI4~\cite{axi}. The core is a host, and a host written
as hardware, so it holds the link's channel ends themselves rather
than the simulation-side \code{Host}: it issues a burst of one beat
per load or store on \code{issue}, sends a store's beat on
\code{wbeat} in the same cycle, and hands the identifier the tracker
granted back on \code{release} as the answer arrives. It writes no
identifier and counts no beat.

Three devices sit behind the link's router, the \code{Router} of
\code{txhdl\_parts}, whose map type states the address map: the data
memory at \code{0x1000}, the serial port at \code{0x3000}, and the
timer and software interrupt the 64~KiB from \code{0x0200\_0000}, and
every other address a hole the router answers \code{DecErr} itself.
The run in \code{run.rs} adds a fourth range, the boot memory at
\code{0x0000}, for the loader. The data memory and the
timer are each behind an \code{AxiPer}, so what each sees is a whole
request and, for a write, its beat. The serial port is an AXI-Lite
peripheral behind \code{LiteBridge<1, ..>}, a bridge of one peripheral,
which sends the router's bursts to it one beat at a time, as
AXI-Lite transactions with no identifier~\cite{axi}. None of the
three decodes an address any more, because the router and the
bridge have already done it. Nothing in this section is the
core's or the devices' own code: the link, the router and the two
trackers are parts, and this design is their first user that is not
a test.

A load waits for its answer and a store is posted, so the core has
one read and as many writes in flight as the link has identifiers,
which is four. A write's response and a load's beat come back on
channels of their own and one identifier goes back a cycle, so a
write's is taken first and a load's waits; that cannot starve the
load, because a load holds the core in its wait and no further store
is issued while it does.

The data memory, Listing~\ref{lst:v-dmem}, is the first device:
1024 words in four memories of a byte, one per lane, so that a store
of a byte or a half writes its lanes and reads nothing. A write lands
in the lanes the request covers; a read goes into the memory's own
register at the edge, the synchronous read a block RAM has, and is
answered the cycle after. The router sends it only the requests in
its range, so it checks no address. What the memory costs as a unit
of its own is the bus's round trip: a load's request crosses the
core's channel, the router and the memory's channel, the memory
reads and answers, and the answer crosses the memory's channel, the
router and the core's channel, so a load retires ten cycles after
its execute, the nine bubbles before each load in the run, where the memory in the core had it retire in one. The
random programs pay it in the wait in writeback and nothing else,
and the board's own memory could take the same place. The
timer, Listing~\ref{lst:v-timer}, is the second device, in the
same subset. It takes every request offered,
answers a read in the cycle after with the word addressed, or zero
for an address that is not its own, writes the lanes a write covers
into its word, counts from the reset one a cycle, and raises its
line while the count has reached the compare; the reset puts the
compare at all ones, since a compare of zero beside a count of zero
is a line raised from the first cycle (issue 419). Its five words are
declared once, with \code{regmap!}, which writes the read mux, the
write enables and the offsets the model and the tests address by; a
word the map does not name reads zero and takes no write (issues 668
and 681). The line is a wire,
not a channel, and a wire between units offers no guarantee about the
step, so the timer drives it from a register, and the run and the
test poll the timer before the core: then the core reads in a cycle
what the timer decided the cycle before, which is what the netlist
does too. The board's other device that raises an interrupt, the
interrupt controller, has a section of its own,
Section~\ref{sec:v-plic}. The timer's own netlist is checked as the core's is, from
the same trace; the two entities share one run. Between two lowered
units a unit's port is one side of the handshake, and the sides wired
together would make a combinational loop, which the placer refused
when the first board top did that. So the board, one lowered unit of
units, joins its children through a channel of the lowering's own:
its netlist puts an instance of \code{txhdl\_chan}, a buffer of two
as the runtime's channel is, between each sender and its receiver.

\lstinputlisting[caption={\code{dmem.rs}: the data memory, an AXI
peripheral.},label={lst:v-dmem}]{cpu/vreteno/src/dmem.rs}

\lstinputlisting[caption={\code{timer.rs}: the timer, an AXI
peripheral.},label={lst:v-timer}]{cpu/vreteno/src/timer.rs}

The serial port, Listing~\ref{lst:v-uart}, is the third device,
with a line each way, and its registers are SiFive's
\code{sifive,uart0} (issue 1011), so that the serial drivers of Linux,
OpenSBI and Zephyr drive it as they are rather than one written here.
It is an AXI-Lite peripheral, so it holds no state for the bus: a read
takes an address and answers the word in the same cycle, and a write
takes an address and a word together and answers at once. A byte
written to \code{txdata} joins a queue of eight, and while the queue
holds any and \code{txctrl}'s enable is set, the oldest goes out on
the line as a start bit, the eight bits least significant first, and
a stop bit, or two, each \code{div} plus one cycles long; a write
while the queue is full is dropped, and a read of \code{txdata} says
so in its bit 31. A frame coming in on the other line is sampled in
the middle of each bit, once a low on the resting line has said a
start bit is there: a high where the start bit should be is a false
start and the frame is dropped, the eight bits in between are shifted
in, and a high at the stop bit lands the byte in a queue of eight
behind \code{rxdata}, whose read takes the oldest and says in bit 31
when there was none. Each queue is a memory of eight bytes with the
index of the oldest and a count, a push and a pop in one cycle leaving
the count where it was; so a terminal may type ahead of the program,
as the terminal of the runs does, its bytes back to back, and a
program may write a line without waiting between its bytes.
\code{ip} says which watermark is passed, the transmit queue holding
fewer than \code{txctrl}'s count or the receive queue more than
\code{rxctrl}'s, and the port's interrupt line is high while one that
\code{ie} enables is; on the board it is source 1 of the interrupt
controller, whose line is the core's external interrupt, and the
demonstration joins it to the line it drives. Two resets differ from
SiFive's, so that a program that never writes the controls finds the
port as it was before the map: both enables start set, and
\code{div} starts at the unit's parameter \code{DIV} less one:
867 for 115\,200 baud at 100~MHz on the board, and 3 in the runs that
are checked, so that a byte takes forty cycles rather than nine
thousand. \code{ie} starts clear, as SiFive's does: a receive bit that
started set let the serial loader's bytes leave a request pending in
the interrupt controller, which Linux took before its port was ready
(issues~\#1136 and~\#1137). The two are the same unit lowered twice, and a driver sets
\code{div} from the clock it is told the port runs on. The line in
goes through one register before the logic, as a pin's should. The
port's pieces, the frame a byte goes out as, the frame after a bit, a
bit's middle and a bit taken in, are functions under \code{\#[lower]},
inlined where the step calls them, so the step is the two sides'
drives; the examples document~\cite{examples} states the form. Its
seven words are declared once, with \code{regmap!}, \code{rxdata} as a
register the read takes, and the macro writes the word a load reads,
which word a write goes to, when a read takes a byte, and the offsets
a test holds to the ones SiFive's drivers use (issue 669). The run and the lockstep
test put a terminal on the two lines, \code{term.rs}: it decodes
what the core sends, and once the core has said a line it types
three bytes back, a pause between them for the program to take
each, and the run prints what was said both ways.
Figures~\ref{fig:v-out}, \ref{fig:v-in} and \ref{fig:v-bus} are
the port at work, three windows on the run's one trace. The first is
the demonstration's O going out: the write reaches the port, the
bit count loads ten, and the line falls for the start bit, then
transmits the byte's bits four cycles each, low first, as the shift
register empties, and rises for the stop bit; the requests after
the first are the program polling the status for the next byte. The second is the
terminal's y coming in: the line falls, the port's register of it
follows a cycle later, the bit count loads ten and counts down as
the bits are shifted in at the middle of each, and at the stop bit
the byte lands in the buffer, whose count rises and with it the
interrupt line.
The third is the bus under the status poll and the write before the
first: the core's request goes out, the router hands it to the
bridge on \code{ar2}, the bridge hands it to the port on \code{uar},
the port answers on \code{ur}, and the answer comes back through the
bridge and the router into the core, which holds in writeback until
it does; then the load retires. The store goes out the same way on
\code{aw2} and \code{uaw}, is answered on \code{ub}, and retires at
once.

\begin{figure}[htbp]
\centering
\input{vreteno_out_timing.tex}
\caption{A byte out: the write into the port, ten bits to send, and
the line's start bit, eight bits low first, and stop bit, four
cycles each; the requests after the first are the program's polls
of the status for the next byte.}
\label{fig:v-out}
\end{figure}

\begin{figure}[htbp]
\centering
\input{vreteno_in_timing.tex}
\caption{A byte in: the line, its register, the bits left, the byte
shifting in, and the buffer's count and the interrupt as it lands.}
\label{fig:v-in}
\end{figure}

\begin{figure}[htbp]
\centering
\input{vreteno_bus_timing.tex}
\caption{The bus under a status poll and a byte write: request out,
through the router and the bridge to the port, response back, the
core waiting in writeback, and the retire.}
\label{fig:v-bus}
\end{figure}

\lstinputlisting[caption={\code{uart.rs}: the serial port, an
AXI-Lite peripheral.},label={lst:v-uart}]{cpu/vreteno/src/uart.rs}

\lstinputlisting[caption={\code{term.rs}: the terminal on the port's
lines, for the run and the test.},label={lst:v-term}]{cpu/vreteno/src/term.rs}

The stages meet at the end of the cycle, in the deferred drives of
the fetch's state and the writeback's. When the execute stage's next
program counter is the one after its own, the word the fetch stage
read this cycle is the right one and moves into the instruction
register, valid, with the fetch counter advanced by two or four, as
the instruction is compressed or not. When it is
not, on a taken branch, a jump, a trap or a return, that word is the
wrong one: the instruction register is marked empty, the fetch
counter takes the target, and the next cycle executes nothing, which
is the one-cycle penalty a machine pays for resolving branches in
its second stage. Reset empties the instruction register and
restarts the fetch at zero; a halting instruction stops the fetch
and parks its counter one instruction past it, as the model's
counter parks, and
sets the halt a cycle later when it retires, so the two agree there
as well. The fetch counter's drive is shaped for the netlist. The
redirect is the last thing the cycle decides, a compare or an add on
a forwarded value. So the counter's next value is two candidates,
advance and target, with reset, park and hold folded into each, and
the redirect chooses between them one multiplexer from the
instruction memory's address. The instruction register's drives are
a \code{case!} on the early conditions alone.

This is the lowerable form, and it reads as a netlist does: every
name after the wait is a function of the state and the inputs, and
nothing branches. \code{select!} chooses among values, the ALU by
\code{f3}, the branch condition, the loaded and the stored word, the
writeback value, the target, the old and the new CSR value and
whether the core knows the opcode; \code{with!} and \code{case!}
choose among drives, the register file's write under its enable,
the control registers under a CSR write, a
trap, the line or a return, the writeback registers under execute's run, and
the instruction register by reset, halt or fetch, in that priority.
Register zero reads as zero rather than being protected from writes, since a mask
on a read is a multiplexer and a guard on every arm is not. The
first form of this file, single-cycle and with Rust \code{match}es,
is in the history; the lockstep test passed each form, which is what
made rewriting it safe.

\lstinputlisting[caption={\code{core.rs}: the three-stage
core.},label={lst:v-core}]{cpu/vreteno/src/core.rs}

\section{The Interrupt Controller}
\label{sec:v-plic}

The core has one external interrupt line, and a machine has more
sources than one. On the board a platform-level interrupt controller,
Listing~\ref{lst:v-plic-gateway} and those after it, stands between
them, at \code{0x0c00\_0000}, where RISC-V machines put it. It is a
part of \code{txhdl\_parts}, \code{Plic<N, EDGE>}, one unit for every
count of sources, whose sources are an array of ports and whose
priorities are an array of registers, \code{Plic1} to \code{Plic8}
naming the counts; it is an AXI-Lite peripheral behind a bridge of
its own. Its words are the
RISC-V PLIC's for two targets, the hart's machine mode and its
supervisor mode: a priority per source from \code{0x4},
three bits, where 0 never interrupts; the pending bits at
\code{0x1000}; the machine target's enable bits at \code{0x2000}, its
threshold at \code{0x20\_0000} and its claim word at
\code{0x20\_0004}; and the supervisor target's at \code{0x2080},
\code{0x20\_1000} and \code{0x20\_1004}, the second context's
offsets. Sources are numbered from 1, and a source's number is its
bit. The two targets share the sources and the pending bits, so a
source enabled for both is served by whichever claims it first; the
supervisor's line is \code{mip.SEIP}'s: the core takes it on its
input \code{seirq}, registered as the machine's line is (issue~\#1094).
A read of \code{mip} or \code{sip} shows SEIP as that line ORed with
the bit software may set, and a set or a clear of \code{mip} modifies
only software's bit, so the line never latches into it, as the
privileged specification requires.

A source's line goes through a gateway first,
Listing~\ref{lst:v-plic-gateway}. A level source asks while its line is
high, and an edge source, whose bit is set in the \code{EDGE}
parameter, asks on a rising edge. A request goes forward, and becomes
pending, only when the source has none in service: service runs from
the claim that took the last request to the complete that ended it.
An edge during a service is held, one deep, and asks on the complete;
a level source still high on the complete asks again at once.
The choice, Listing~\ref{lst:v-plic-choice}, is a loop over the
sources that keeps the best so far. A source is a candidate when it is pending,
enabled and above the threshold, and the best has the highest
priority, the lowest number winning a tie. A read of the claim word
answers with the best source's number, or 0, and clears its pending
bit; a write of the number back is the complete. Each target chooses
alike, from its own enables and above its own threshold, and each
target's interrupt line is a register, high a cycle after it has a
best source, so the core reads the machine one as it reads the
timer's.

\code{ex\_plic}, Listing~\ref{lst:v-plic-ex}, drives a \code{Plic3}
whose second source asks on an edge, the way a handler would: it sets
the priorities to 1, 3 and 3 and raises all three lines. The claims
come out 2, 3 and 1, and a fourth claim answers 0. On the completes the
two level sources still high ask again and the edge source does not. A
threshold of 1 hides a priority of 1, and a second pulse on the edge
source during its service asks once on its complete. Last, source 3
is enabled for the supervisor target alone: it asks on that target's
line and not the machine's, and is claimed and completed at the
supervisor's offsets.
Figure~\ref{fig:v-plic} is the start of that run. The build simulates
the controller's netlist against the run under nvc and Verilator, and
\code{//lib/parts:parts\_test} checks the same four rules on their
own.

\begin{figure}[htbp]
\centering
\input{plic_timing.tex}
\caption{Three sources asking at once: the lines, the pending and
active bits, the claims on \code{ar} and their answers on \code{r},
and the interrupt line, which falls a cycle after the last claim and
rises again when source 1, completed and still high, asks again.}
\label{fig:v-plic}
\end{figure}

\plicpart{gateway}{The gateway: the lines as one word, and the
requests that go forward.}{lst:v-plic-gateway}

\plicpart{choice}{The choice of the source a claim takes.}%
{lst:v-plic-choice}

\plicpart{bus}{The words on the bus, the claim and the complete.}%
{lst:v-plic-bus}

\plicpart{drives}{The drives: the requests, the registers and the
line.}{lst:v-plic-drives}

\lstinputlisting[caption={\code{ex\_plic.rs}: a controller of three
sources, claimed and completed.},label={lst:v-plic-ex}]{lib/examples/ex_plic.rs}

\lstinputlisting[style=ascii,caption={What \code{ex\_plic} printed when
the build ran it.},label={lst:v-plic-out}]{out_plic.txt}

On the board the controller is a \code{Plic3}. Its source 1 is the
serial port's interrupt line, its source 2 the board's \code{irq}
input and its source 3 the Ethernet peripheral's, all level sources, and its line is the core's external
interrupt. \code{irq.rs}, Listing~\ref{lst:v-irq}, is a program for the
board that takes its input by interrupt. It enables source 1, installs
the trap entry of \code{vreteno\_hal} (Section~\ref{sec:v-hal}),
registers a handler for the external interrupt's line, allows that
interrupt and says \code{ready}, then waits in a loop that never
reads the serial port. Each byte the terminal types lands in the
port's buffer and raises its line. The controller raises the core's,
and the handler claims, reads one byte and completes. A byte still in
the buffer keeps the port's line high, so the controller asks again on
the complete, and the next byte is the next interrupt. The core's
pending bit for the external interrupt is the controller's line, so
the crate's dispatcher clears nothing after the handler; the
controller's line fell a cycle after the claim. \code{//cpu/vreteno:board\_test} runs the board's design on it,
with a terminal that types \code{ping}, and the program says
\code{got ping in 4 interrupts}.

\lstinputlisting[caption={\code{cpu/vreteno/rust/irq.rs}: input by
interrupt, one interrupt a byte.},label={lst:v-irq}]{cpu/vreteno/rust/irq.rs}

\section{The Programs}
\label{sec:v-programs}

Listing~\ref{lst:v-program} holds the programs. \code{Asm} is the
smallest assembler that has labels: a branch to a label not yet
placed is emitted as a placeholder and patched when the label is,
which is all a forward branch needs. It holds halfwords, a thirty-two
bit instruction being two of them and a compressed one one, and packs
them into the little-endian words the instruction memory holds. With
\code{compress} set, \code{emit} writes an instruction in its
compressed spelling where it has one; a branch to a label keeps
thirty-two bits, since its size is fixed before the label is placed.
\code{align} pads to a whole word, which a trap handler needs:
\code{mtvec} keeps no low bits, so a handler starting in the upper
half of a word is entered two bytes early. The demonstration program
is written with \code{compress} set, so it runs the same
instructions in the same cycles in fewer bytes. It sums
one to ten in a loop, calls a function that doubles the sum and
returns through \code{jalr}, stores and loads every width so the sign
extension of \code{lb} and \code{lh} and the zero extension of
\code{lbu} and \code{lhu} are exercised, uses the upper immediates,
the arithmetic shift and both compares, then traps twice and ends
in the write to \code{mhalt} with 110 in \code{x10} and at the first
data word.
The traps need a label's absolute address, which \code{abs} gives
an encoder once the label is placed. The handler's address goes into
\code{mtvec} first of all, the timer's compare is set to 150, both
interrupts are enabled in \code{mie} and interrupts in
\code{mstatus}, so that the line can be raised at any point of the
run and the timer fires in the middle of it. An \code{ecall} and then the
word of all ones, which is illegal, enter the handler. The handler
reads the cause into \code{x23} and branches on its sign: an
interrupt has nothing to clear in \code{mip}, whose bits are the
lines, and the timer's has the compare
moved past the end of the run and a \code{fence} after it, since the
store is posted and the line drops only when it lands, and either is counted in \code{x8}
and returned from to the word it interrupted; an exception has
\code{mepc} read, stepped past the word that trapped and written
back. The program then writes \code{mscratch} by immediate and
reads it back into \code{x24}, so that two is in \code{x23} and
five in \code{x24} at the end, and two in \code{x8}: the line,
raised once, and the timer. Before the traps it writes OK and a
newline to the serial port, a byte at a time. Each byte waits for a
load of \code{txdata} to say the queue has room, its sign bit clear,
the polling loop every program writes for a port like this. Then it
echoes three bytes: each loads \code{rxdata} until its sign bit says a
byte came, which the load takes, waits for room and writes it back; the
terminal typed the three back to back, so the second and third
waited in the port's buffer. The
timer's interrupt is enabled from the start and the line's only
after the echo, since a byte received raises the line too; the
pulse and the bytes have left the pending bit set by then, so the
line's interrupt comes once, as counted. Between the shifts
and the traps it multiplies and divides, five of the eight
instructions, on the sum and on minus two: minus 220 into
\code{x26}, minus 55 into \code{x29}, and the remainder minus two
into \code{x30}.

The random generator writes straight-line programs of a given length
from a seed: every register operation, aligned stores and loads
within the first data words through \code{x2}, which it keeps at the
data base, now and then a forward branch over one or two words, so
both outcomes of every branch condition occur, the M extension in
half the places an \code{or} would go, the A extension's word forms
on an aligned data word through \code{x28}, an AMO, or \code{lr.w}
and \code{sc.w} on one word, or \code{sc.w} alone; the mode the
body runs in, machine, supervisor or user, with illegal instructions,
user mode's environment calls and the supervisor's software interrupt
each delegated or not by the seed, and a supervisor's handler beside
the machine's; and the six CSR forms on
\code{mscratch}, \code{mcause}, \code{mepc} and \code{mtval},
and now and then an \code{ecall} or an illegal instruction: a word of
zeros, which is two illegal halfwords, a reserved compressed halfword,
or a word of ones. Half of its instructions are shaped to have a
compressed spelling, with the destination as the first operand, the
registers \code{x8} to \code{x15} and a six-bit immediate, and are
written in it, and some forward branches are \code{c.beqz},
\code{c.bnez} or \code{c.j}. So the two lengths are mixed, and a
thirty-two bit instruction often starts in the upper half of a word.
Half the word stores and loads go to the timer through
\code{x30}, which the program keeps at its address, so the count
and the compare are written with anything, and the count read
back. Both interrupts are enabled from the start, and the compare is
zero then, so the timer fires at once and keeps firing. The handler
placed after the halt branches on the cause's sign as the
demonstration's does: an interrupt has the line's bit cleared and
the compare set to the count plus sixty-four, both halves, since a
program may have stored anything into the count, and returns to the
word it interrupted; an exception returns past the instruction that
trapped, four bytes on for an \code{ecall} or an illegal instruction
whose low two bits are both set and two for the rest, which the
handler reads off the cause and the trap value. \code{x30} and
\code{x31} are the handler's and the program leaves them alone;
\code{x29} is the handler's too, and a program that writes it loses
it at the next trap, in the core and the model alike. Nothing in it can loop, so every program halts, the
timer's storms included.

\lstinputlisting[caption={\code{program.rs}: the assembler, the
demonstration, the random
programs.},label={lst:v-program}]{cpu/vreteno/src/program.rs}

\section{The Lockstep Test}
\label{sec:v-lockstep}

Listing~\ref{lst:v-lockstep} is the test, \code{bazel test
//cpu/vreteno:lockstep\_test}. It builds the core with the program, keeps
handles on its program counters, valid bit, register file, data
memory and halt bit, which \code{Reg} and \code{Mem} give out as
clones, holds the reset for one cycle, and then runs the core a cycle
at a time, stepping the model once whenever the writeback says an
instruction retired. After each cycle the architectural program
counter, that of the instruction about to execute or the fetch's
when the stage is empty, the thirty-one registers, the eight control
registers, the two debug registers, the debug mode and the halt must
agree, and the failure names the cycle
and the instruction; at the halt every word of data memory must
agree too, and the count of
retired instructions must not exceed the count of cycles, which is
what one instruction per cycle at most means. The demonstration program passes with the values the
comments promise, and sixty-four random programs of two hundred
words each pass, more than two hundred instructions where some are
compressed, which covers every instruction the generator
emits many times over with operands it did not choose. The test also
walks each random program from its start and counts its instructions
by length: across the sixty-four there must be more than three
thousand compressed ones, more than three thousand whole words, and
more than two thousand thirty-two bit instructions that start in the
upper half of a word, so the mix the fetch has to get right cannot
quietly go away.

The compressed instructions have a test of their own,
\code{compressed\_instructions}, which runs every one of them at
least once with the lengths mixed: the arithmetic group, the stack
pointer's short forms, the loads and stores, a loop on
\code{c.bnez}, calls by \code{c.jal}, \code{jal} and \code{c.jalr}
with their links two or four bytes on, returns by \code{c.jr}, a
reserved halfword that traps with itself as the trap value, and the
write to \code{mhalt}, a thirty-two bit instruction among the short
ones, to end. Two more negative controls belong to it. With
a compressed call linking four bytes on, it fails at cycle 59 on
\code{x1}; with a thirty-two bit instruction in the upper half of a
word taking its upper half from its own word rather than the next, it
fails at cycle 9 on the first such instruction.

A test that cannot fail proves nothing, so negative controls were
run, six of them. With the arithmetic shift in the ALU replaced by the
logical one, the demonstration fails at its \code{srai}, naming the
register, the cycle and the instruction, and the random programs
fail within a few cycles of their first \code{sra}. With the redirect
kept but the squash removed, so that the word after a taken branch
executes, the demonstration fails at its first \code{bne}. With the
forwarding path cut, it fails at the first instruction that reads
what the one before it wrote. And with \code{mret} restoring the
interrupt enable but leaving the previous-enable bit clear, where
the manual sets it, the demonstration fails at its first
\code{mret}, on \code{mstatus}. With \code{mtval} written as
zero for an illegal word, the demonstration fails on \code{mtval}
at that word, which is why the word is all ones rather than zero.
With the timer's compare made strict, so that it fires a cycle late,
a random program fails on the program counter. Every control was
then removed. The single-cycle core passed the
first; the pipeline passes them all. The random programs found a
defect the controls did not look for: a division interrupted in the
middle went on stepping under the handler's instructions, and one of
them was a multiply, so the pair was overwritten and the quotient
came back wrong nineteen steps short. The sequencer is cancelled by
an interrupt now.

The line is the test's. For the demonstration program it is
raised at cycle forty and held until the core takes it, since the
pending bit is the line and not a latch; for a random program it is high in one cycle of sixteen
from the seed. The core decides on the pending bit as it stood
before the cycle, so the test reads that bit with the two enables
while the instruction is in execute, and tells the model, when the
instruction retires a cycle later, whether the interrupt was taken
in its place; the model raises its own pending bit in the cycle the
line is high, as the core does at its edge. The timer's line and
the bus's answers go across the same way. The test reads the timer's
pending register while the instruction is in execute, and the core's
register for a load's answer before the cycle in which the load
retires, and hands both to the model with the instruction. So a load
from a device agrees on what the device said, and the timer's
interrupt is decided on the same bit in both. The test hands the model
the interrupt the core decided on, a cycle old, and checks each one
against the model's enables where it is taken, which is what fails
if the decision is not dropped after a write of them (issue 1331). The compare lands in
the timer a cycle after the core's store, so it is checked at the
halt, as the memory is. So are the serial port's count of bytes and
its last byte, against the bytes the model was given.

The control registers are compared with one allowance. The core
writes them in execute, when the instruction is in its second stage;
the model writes them when the instruction retires, a cycle later.
So in the one cycle in which a CSR instruction, a call, a return,
an illegal word, an interrupt or a store, which may be to the
timer's compare, is in execute the two disagree by design; the
compare is checked with the control registers. The test
skips the comparison in that cycle and requires it in every other,
and the negative control above shows the allowance is narrow enough
to catch a wrong status.

The debugger is the test's as well, as a plan: the cycle of a halt
request, held until the core is in debug mode; the cycles to hold it
there; and how many resumes to give, each two cycles high, since the
core acts on the level once per entry and the line must drop between.
The model has no debugger, so it follows the core's mode register at
each edge: entered, it records the cause and \code{dpc} as the core
does, before the instruction in execute and after the one behind it
retired; left, it goes to \code{dpc} and arms a step when
\code{dcsr.step} is set. In debug mode the architectural program
counter is \code{dpc}, since nothing is in flight and the fetch has
run ahead. \code{dpc} and \code{dcsr} are compared with the other
control registers and the mode itself in every cycle, so a core that
ran an instruction while halted, or entered a cycle late, fails on
the program counter or on a register. Three programs use the plan,
which also names what the debugger writes through the module's port
at an entry, as OpenOCD writes \code{dcsr}.
The demonstration, halted at cycle two hundred, held for twenty
cycles and resumed, ends as it does without the debugger, with one
entry whose cause is the request. A program halted once, with
\code{dcsr.step} written at that entry and resumed as often as the
plan asks, enters again after every instruction with the step as the
cause until the tenth entry clears the bit, and then runs to its halt
with every step counted in a register. And a program halted at its
start, with \code{dcsr.ebreakm} written and resumed once, enters
at its \code{ebreak}, with the breakpoint as the cause and the
instruction's address in \code{dpc}, and is left there, held. A
fourth program, with no plan, reads and writes both registers from
machine mode and traps on every access. The plan with no resume left ends the run in the
mode, which is the one run that does not end in the halt.

The reset line is the test's too, high for one cycle when a run asks
for it. A program that has set its trap vector, enabled the timer's
interrupt and enabled interrupts is reset in its loop: the core starts
again at zero with the CSRs cleared and the timer's compare all ones,
the model is reset the same way once whatever retired in that cycle
has been stepped, and the two agree in every cycle after it. The
register file keeps what it held, so the program's count of its own
starts is two the second time, and it halts on that.

The bus's refusals go across to the model as the timer's line does:
the core's register for a refused load is read with the load's
answer, and the model traps on the load instead of writing; the
register for a refused store is read before each cycle, and when it
is set the model is told the fault was taken before the instruction
that retires next, as it is told of an interrupt, so \code{mepc}
agrees however long the bus took. The program for it loads from and
stores to an address nothing decodes, with a handler that counts,
sums the causes, keeps the trap values and steps past a load, then
sets \code{mbusquiet} and does both again to read a zero and drop the
store. The board test's program that pokes the boot memory has a
handler for the refusal too. A third program jumps into nowhere and
comes back through its handler with the fetch's cause and address.
The fence has a negative control of its own: with the instruction
retiring at once, as it did before issue 432, the demonstration's
handler returns before its compare has landed and the run counts
three interrupts where it counts two.

\lstinputlisting[caption={\code{lockstep.rs}: the core against the
model.},label={lst:v-lockstep}]{cpu/vreteno/tests/lockstep.rs}

\section{The Run}
\label{sec:v-run}

Listing~\ref{lst:v-main} runs the demonstration program and prints a
line per cycle, the program counter, the disassembly of the word
executed and the register written, or a bubble when the execute
stage was empty; Listing~\ref{lst:v-out} is what it printed when the
build ran it. An instruction the program holds compressed is marked
with a \code{c} and printed as the instruction it stands for, which
is what the core executes; the address after it is two bytes on
rather than four. The loop is visible as the three
instructions at \code{0x36} that repeat ten times with a bubble after
each taken \code{bne}, the call as the jump to \code{0x110} and the
compressed return to \code{0x42}, each with its bubble, and the loads
as the sign and zero extensions of the same byte and half, each after
the nine bubbles of the bus's round trip. The one stall a load causes
after it is the bubble after \code{lw}, where the next instruction
adds to what it loaded; each multiply is a run of two bubbles and
each divide a run of thirty-three, which the run prints as one line
with the count, and the last line of the printout is what the port
said, decoded from its line. The timer's interrupt is the first jump
to the handler at \code{0x114}, at \code{t} 308 in place of the first
\code{div}, once the timer reaches the compare of 150: cause seven
and the top bit in \code{x23}, the compare moved past the end of the
run, one counted into \code{x8}, and the return to the \code{div} it
displaced, which starts again. The line's interrupt comes next. The
line is raised at cycle forty and held until the core takes it, as a
device holds its line until it is served, since the core's pending
bit is the line rather than a latch (issue 958). The moment the echo's
\code{csrrw} at \code{0x0f8} enables it, at \code{t} 1140, it is taken
in place of the \code{ecall} after it: cause eleven and the top bit in
\code{x23}, two counted into \code{x8}, and the return to the
\code{ecall}, which traps in its turn. The two traps are the
last jumps to the handler, the first from the \code{ecall} with
eleven read into \code{x23}, the second from the illegal word, all
ones and thirty-two bits long, with two, each with its bubble, as the
returns have theirs. Figure~\ref{fig:v-wave} is the opening of the same run, the reset
and the first dozen instructions, drawn from the waveform the run
wrote; \code{pc} there is the fetch
stage's counter, one instruction ahead of the one executing, and
\code{done} is low on the bubble. The whole run is in the FST beside
it, for a viewer. The run ends by writing the VHDL of the core, with
the program in its instruction memory, which Section~\ref{sec:v-netlist}
takes up, and by printing the Verilog, which closes
Listing~\ref{lst:v-out}.

\lstinputlisting[caption={\code{main.rs}: the run and the
waveform.},label={lst:v-main}]{cpu/vreteno/src/main.rs}

\begin{figure}[htbp]
\centering
\input{vreteno_timing.tex}
\caption{The opening cycles: reset, then an instruction per edge,
the handler's address, the timer's compare and the two enables set
up, the counter moving on by two after a compressed instruction and
by four after the rest; \code{rd} and \code{val} are the writeback,
\code{rd} zero when nothing was written.}
\label{fig:v-wave}
\end{figure}

\lstinputlisting[style=ascii,caption={What \code{vreteno} printed when
the build ran it.},label={lst:v-out}]{out_vreteno.txt}

\section{The Interrupt Controller}
\label{sec:v-clint}

The timer is the core's interrupt controller, placed where
standard RISC-V platforms locate one: it fulfills the standard
requirements that an operating system port asks of a platform, providing
\code{mtime} and \code{mtimecmp} at the conventional offsets, and
\code{msip}, one bit whose writing raises a software interrupt.

\begin{center}\footnotesize
\begin{tabular}{@{}lll@{}}
\toprule
Offset & Name & What it is \\
\midrule
\code{0x0000} & \code{msip} & bit 0 raises the software interrupt \\
\code{0x4000} & \code{mtimecmp} & the compare, low half \\
\code{0x4004} & & its high half \\
\code{0xbff8} & \code{mtime} & the count, low half \\
\code{0xbffc} & & its high half \\
\bottomrule
\end{tabular}
\end{center}

The window is 64~KiB, which those offsets need, so the controller has
a range of its own in the address map rather than the nibble the timer
had. The offsets are far apart, so the decode is on the offset and not
on two bits of it, and a word in the window that is none of the five
reads as zero.

\code{msip} is a register and a line: a program writes one to raise
its own interrupt and zero to clear it, and the line is the register
itself, so nothing in between can lose it. That is how a scheduler is
entered: a port raises the software interrupt, the trap takes the
machine into the handler, and the handler clears the bit and switches
context.

The core gains a line for it, \code{mip} bit 3 beside the timer's bit
7 and the external bit 11, and \code{mie} gates it as it gates the
others. When more than one is pending the core takes the external
interrupt first, then the software one, then the timer's, which is the
order the specification gives and the order the reference model keeps,
so the lockstep compares them.

\code{//cpu/vreteno:lockstep\_test} runs a program that raises the
software interrupt for itself three times and counts the traps in the
handler, against the model, cycle by cycle.

\section{Fetching From the Bus}
\label{sec:v-fetch}

The core fetched only from the memory inside it, so the program was
fixed when the bitstream was, and four kilobytes was all of it. That
was the one thing in the way of everything else: a bootloader, an
operating system that starts a task from an image, a test runner on
the target.

The counter now says where the instruction comes from. Below
\code{IMEM\_BYTES} it is the memory inside the core, as before, and
that memory is the boot ROM: what runs before there is anything else.
At or above it, the core asks the bus, on the same channels its loads
and stores use, and waits.

The same words sit on the bus as well, at address zero, behind a
tracker like the data memory's, readable and not writable
(\code{crate::rom}, issue 268). The fetch reads the copy inside the
core in the same cycle; a load reads the copy on the bus a few cycles
later, like any memory; the netlist initialises both from one image,
so they cannot disagree. A write to it is taken, its beat consumed so
the channel does not jam, and answered \code{SlvErr}, which is what
settles what a fetch sees when a write lands on the same address in
the same cycle: nothing lands. So a program's constants live beside
its code, \code{.rodata} is linked into the boot memory, and a
debugger on the bus can read the program that is there. Every load
and store goes to the bus; the core no longer keeps a hole below the
data memory.

The primary latency source is sequential transaction stalling.
The core issues one transaction at a time, so an instruction fetch goes
out only after outstanding data transactions complete, and the pipeline
stalls until the requested word returns. A thirty-two bit
instruction at an odd halfword takes its upper half from the word
after, so the fetch keeps a buffer of two words and asks for the
second only once the first says it is wanted; a run straight through
memory therefore asks for a word every second instruction rather than
every one.

It was correct and it was slow: issue~\#1009's loop of three
instructions ran at 14.655 cycles an instruction from the data memory
and at 1.333 from the boot memory. So a fetch from the data memory or
the DDR3, the two memories behind the bus, goes to an instruction
cache instead (issue~\#1021): sixteen kilobytes, 1024 lines of four
words, one way. The fetch registers the word's physical address where
it would have sent it to the bus, once translated, and the cycle after
reads the line's tag and the word at it; so the cache is indexed and
tagged by the physical address, and two virtual pages of one physical
page share its lines. It was four kilobytes at first, and the GL
frame's code, of 34~KB, did not stay in it: at sixteen its misses
fell from 6.45 to 0.70 per cent of lookups and the frame took 25 per
cent fewer cycles; thirty-two bought three more (issue~\#1290). Its
words are in two banks, the even and the odd: a lookup reads a word
and the one after it, which are always in different banks, so each
bank is read once a cycle and is a block RAM (issue~\#1303). A hit puts the
word in the buffer; a miss asks the bus for the line as one burst of
four, writes each beat into it, and looks it up again. The data memory
answers such a burst a beat a cycle. That brought the loop from the
data memory to 4.093 cycles an instruction, most of what was left
being the buffer, which served one instruction's words and not the
next one's, so each instruction was a lookup of its own.

So a hit fills both of the buffer's words when the line holds the
next one, and the word after the buffer's one is asked for ahead of
time when it is in the same line (issue~\#1187). An instruction is
read from the first word only. When the counter moves on from an
instruction there to one that begins in the second, the buffer shifts
along at that edge, the second word and its fault marks moving to the
first; if the second word is still out, it goes into the first place
when it comes back. A jump into the second word misses and fetches it
again, which is rare. The first version compared the counter with
both words' addresses and chose between them before the instruction
register, and the flagship missed timing on that choice
(issue~\#1279); shifting at the edge keeps one compare, of a register
with the counter, in front of the read and the fetch.
\code{//cpu/vreteno:icache\_test} measures the loop at 3.349 cycles an
instruction. It counts from the reset, and the loop jumps to the data
memory at once, so it waits first for the tags to be cleared: 1024
cycles at sixteen kilobytes, 0.341 of a cycle over its 3002
instructions. At four kilobytes the clearing took 256 cycles, and the
test measured 3.093 with the window. Asking ahead alone moved
nothing, since one fetch is out at a time and a lookup takes three
cycles: two words from one lookup is what pays. A fetch can now be
out under an instruction that is ready to run, so a load waits while
one is: the core still holds one read at a time.

A store does not reach the cache, as the specification allows, so code
written by stores runs after a \code{fence.i}, which clears the tags,
a line a cycle, as a reset does. The loader in the boot memory ends
with one before it jumps to the program it was sent.

\code{//cpu/vreteno:lockstep\_test} runs a program that lives above the
boot memory: the boot memory holds a jump into the data memory, the
program itself is in the data memory, and every instruction after the
jump comes back over the bus. The model fetches from the same place
and the two agree cycle by cycle, which is what says the fetch is a
fetch and not a guess.

\section{Virtual Memory}
\label{sec:v-mmu}

Linux runs each program in an address space of its own, and the
hardware that makes that possible is the memory management unit. On a
thirty-two bit RISC-V hart it is Sv32: a virtual address is translated
through a table of two levels into a physical one, page by page, and
an access the table does not allow traps (issue~\#1014). This section
is the unit, \code{Mmu} in \code{lib/parts/src/mmu.rs}, and
Section~\ref{sec:v-hart} the core around it. Where it goes was decided
by measuring first,
in \code{docs/sv32-timing.md} (issue~\#1009): the core asks with an
address it already holds in a register, so the lookup is off the two
paths that set the clock, at the cost of a cycle per access that goes
to the bus.

The unit has two translation lookaside buffers, one for fetches and one
for data, each of eight entries, fully associative and refilled in
turn, and one walker that fills both. Listing~\ref{lst:v-mmu-lookup}
is the lookup. Every entry is compared with the request's page at
once; a megapage, a leaf at the first level, covers four megabytes and
compares only the upper ten bits. The answer is a register, so it
comes a cycle after the request: on a hit, the physical address, or a
page fault if the entry does not allow the access. A miss answers
nothing, and the core holds its request, the address and the request
bit both steady, while the walker reads the table; the answer comes a
cycle after the entry is filled. An answer always belongs to the
request of the cycle before, so one that arrives after the core has
moved on is for the old request, and the core ignores it. Outside
supervisor and user mode, or with \code{satp.MODE} clear, the address
passes through unchanged, with the same latency of a cycle, so the
core has one rule for every access.

The check, Listing~\ref{lst:v-mmu-check}, is the specification's:
execute for a fetch, write for a store, read for a load or execute
under \code{MXR}; a user page only from user mode, and from supervisor
mode only for data under \code{SUM}; and the accessed bit always, with
the dirty bit for a store. The walker, Listing~\ref{lst:v-mmu-walk},
never writes. An entry whose accessed or dirty bit is clear is a page
fault, and the kernel sets the bit and flushes, which the
specification allows; a walker that wrote would need an atomic read
and write on the bus. It reads one entry at a time over a channel of
its own, which the core will put on its bus port: \code{satp}'s page
first, indexed by the upper ten bits of the page number, then the
second level's if the first entry points to one. An invalid entry, a
pointer at the second level and a megapage whose low page number is
not zero are page faults; a read that fails, or a page past the
board's thirty-two bits of physical address, is an access fault.
A flush, on a write of \code{satp} or an \code{sfence.vma}, empties
both buffers, and a walk under way when it comes finishes without
filling anything. A request in the cycle of a flush is not answered
from the buffers as they were: it is held, and answered from the
emptied ones by a walk. Address space identifiers are not kept, which the
specification allows.

\mmupart{lookup}{The lookup: every entry at once, and the answer in the
cycle.}{lst:v-mmu-lookup}

\mmupart{check}{The check an entry's bits must pass.}{lst:v-mmu-check}

\mmupart{walk}{The walker: the first level's entry, the answer, and
what makes it a fault.}{lst:v-mmu-walk}

\code{translate} in the same file is the reference: the
specification's walk with no buffer at all. The tests in
\code{lib/parts/src/mmu/tests.rs} hold the unit to it, against tables
in a memory that answers the walker after a few cycles: a page walked
once and then hit a cycle after it is asked, a megapage, every fault,
the flush, a flush in the middle of a walk and a request in the cycle
of one, the turn in which entries are replaced, and six hundred random
requests over random tables, on both ports and in every mode.

\code{//lib/parts:mmu\_pnr} places and routes \code{Mmu8} alone, out of
context, against the core's ten nanosecond clock, with the requests
arriving a nanosecond after the edge, as from a register of the core's.
It meets the clock with 2.440\,ns to spare, in 493 LUTs and 919
flip-flops, the eight entries of each buffer held in flip-flops rather
than block RAM. The lookup, from the request's address through every
entry's compare and the check into the answer's registers, has
3.583\,ns to spare for a fetch and 3.713\,ns for data, and the answers
leave those registers 1.22\,ns before they reach the unit's edge. The
answer was not always a register. Used in the cycle it was asked, it
took 7.86\,ns from the address to the physical address, and the core's
bus request already spends 8.9\,ns of its own
(\code{docs/sv32-timing.md}), so the two would not have fit in one
cycle. The register costs the one cycle per access that the note
already counts, and the board's place and route will say what the
paths are in place.

\code{ex\_mmu}, Listing~\ref{lst:v-mmu-ex}, is the unit as a kernel
meets it. A fetch from the kernel's text walks one level to a megapage
and then hits; a load from a user page faults until \code{SUM} is set;
a store to a page whose dirty bit is clear faults until the kernel sets
it and flushes; and with both ports missing at once, the data port's
walk goes first, since its instruction is the older.
Figure~\ref{fig:v-mmu} is the first walk. The build simulates the
unit's netlist against the run under nvc and Verilator.

\begin{figure}[htbp]
\centering
\input{mmu_timing.tex}
\caption{A fetch from the kernel's text: the request, the walk's one
read on \code{ptw} and its answer on \code{pte}, and the answer to the
fetch, then the next fetch from the same megapage, a hit answered a
cycle after it is asked.}
\label{fig:v-mmu}
\end{figure}

\lstinputlisting[caption={\code{ex\_mmu.rs}: translation as a kernel
meets it.},label={lst:v-mmu-ex}]{lib/examples/ex_mmu.rs}

\lstinputlisting[style=ascii,caption={What \code{ex\_mmu} printed when
the build ran it.},label={lst:v-mmu-out}]{out_mmu.txt}

\subsection{In the core}
\label{sec:v-hart}

\code{Hart}, Listing~\ref{lst:v-hart}, joins the unit and the core
with the core's ports and nothing more, so the board, the pair, the
runs and the tests hold a hart where they held a core. The unit goes
first in the join, since its answers are registers the core reads in
the same step. The core's requests are registers too, and the unit
reads them a step late in simulation and on time in hardware; so the
core believes an answer only two cycles after it made the request,
which is right in both, since the unit holds its answer as long as the
request is held. The walker's reads go out through the core's own port
on the bus, one at a time, once every store the core has posted has
been answered. Only one read of the fetch's, the data's or the
walker's is ever out.

Translation is on below machine mode when \code{satp} says so. The
fetch then takes every word from the bus, the boot memory's included,
since a page may map anywhere. It keeps one translation, of the page it
last asked for, in \code{ft\_valid}, \code{ft\_vpn} and \code{ft\_ppn},
and asks the unit only when the word it wants is on another page, so a
fetch within a page costs nothing more. A page that faults is kept the
same way, and its words are taken as faults at once, without a read.
The fault is the first half's if that half's page faults, else the
second's, which is on the next page only for a thirty-two bit
instruction in a page's last halfword; the trap's value is the address
of the half that faulted.

A load, a store, \code{lr.w}, \code{sc.w} or an AMO waits in execute
for its physical address. The address the adder computed goes into
\code{x\_va}, the unit is asked from that register, and the answer is
kept in \code{x\_pa}, \code{x\_pf} and \code{x\_af} until the
instruction leaves execute. A page fault is cause 13 for a load and
\code{lr.w} and 15 for the rest, which write; a walk that reads nothing
is an access fault, 5 or 7. Either is taken from those registers, so
neither the choice of trap nor the redirect waits on anything late. A
misaligned access traps as one whatever its page says, as the
specification orders the two. In machine mode with \code{MPRV} set, a
load or a store is translated and checked as the mode \code{MPP} names,
and the fetch is not; that is how OpenSBI reads the instruction a
supervisor trapped on, which Linux makes it do before its first
program runs. The core tells the unit the data's mode, which is the
fetch's but under \code{MPRV}, and a return below machine mode clears
the bit (issue~\#1105).

A write of \code{satp} and an \code{sfence.vma} tell the unit to flush,
and refetch the instruction after them, since that one was fetched
under the old translations. \code{sfence.vma} first waits, as a fence
does, until every store is answered, so that the tables a kernel wrote
are in memory before a walk reads them. A change of mode with
translation on, a trap, a return or a resume, drops the fetch's one
translation and its buffer, and a fetch out on the bus across any of
these is dropped when it answers. \code{fence.i} does the same to the
buffer: it used not to, and a program that rewrote the word the buffer
held ran the old one (issue~\#1096). It also clears the instruction
cache's tags (issue~\#1021).

The lockstep test holds the hart to the model, which translates with
\code{translate} and keeps nothing. One program writes tables into the
data memory's one page, the root and the second level both, and runs
supervisor mode through every kind of page: allowed and not, under
\code{SUM} and \code{MXR}, a walk that reads nothing, ten pages twice
past a buffer's eight entries, an entry rewritten and flushed, a jump
where no page is, and an instruction that straddles into a page that is
not there; it checks every trap's cause and value. \code{random\_vm}
then runs the random programs' body in supervisor or user mode, over
tables the seed chooses, with the page faults delegated or not, and
machine mode either stepping past one or doing what a kernel does:
setting what the entry lacks, flushing, and running the instruction
again. Ninety-six seeds run in lockstep. The cost is the fetch: under
translation every word comes from the bus, and every return translates
its page again; the instruction cache of issue~\#1021 takes the words
from its lines, and a return still translates its page.

The flagship, Vreteno's board with the Ethernet and HDMI parts on one
FPGA, places and routes with the hart in it. On 2026-10-04 at commit
\code{f9be2042} it met the core's ten nanoseconds with 0.099~ns of
slack, and hold with 0.028~ns, in 21\,327 LUTs and 19\,786
flip-flops, as much slack as it had with user and supervisor mode
alone (issue~\#1012). Its worst path is now in the
fetch's buffer, from the address the buffer holds through the compare
with the one the fetch wants and the add of the next word's, into the
buffer's register again. It took two tries. The first run missed by
1.599~ns: the delegation of a trap, which chooses between \code{stvec}
and \code{mtvec} and so feeds the next program counter, had the new
faults ahead of the late choices, so the stall's settling went through
two more multiplexers on its way to the boot memory's address. With the
faults chosen first, from registers, it missed by 0.274~ns, on the
stall itself, which waited on the adder twice, once for the bus and
once for the translation. A misaligned access now waits as an aligned
one does and traps when it runs, so the stall waits on no adder at all,
and the path that had set the clock since the stall was written is no
longer the worst.

\lstinputlisting[caption={\code{hart.rs}: the core and its unit, with
the core's ports.},label={lst:v-hart}]{cpu/vreteno/src/hart.rs}

\section{The Netlist}
\label{sec:v-netlist}

The Verilog at the end of Listing~\ref{lst:v-out} is the core as
hardware: its registers, two memories, the instruction
memory holding the program, a wire per \code{let} of the loop, seven
temporaries, and one clocked block of a hundred and twelve lines;
the timer's, after it, is six registers and a block of twenty-eight.
The VHDL the run
wrote beside it says the same, and the build simulates both against
the trace the Rust core wrote, the VHDL under nvc and the Verilog
under Verilator, the testbench replaying the reset and checking every
register and every output at every tick; the netlist agrees. The
tests are \code{//docs:vreteno\_sim\_vreteno\_vreteno\_tb\_test}
and \code{//docs:vreteno\_vsim\_vreteno\_test}: the run writes the
netlists of nine units, and six of them are checked against the run
this way, so each name designates the entity as well as the run's.

The lowering grew in five ways for this, each stated and checked in
the examples document~\cite{examples} before the core used it. A
memory is an array, written through \code{at} and read by index, and
the netlist takes the memory's first words from the example, since
the lowering cannot see what \code{Mem::with} gave at run time.
\code{slice}, \code{concat}, \code{sext} and \code{zext} are the
part selects, concatenations and extensions of the target language,
the width of an expression inferred where an extension replicates a
bit; a slice of an expression, which neither language allows, is
hoisted into a wire. \code{select!} is a chain of conditionals, and
an integer is a pattern. A \code{let} is a wire named for it, so the
netlist reads as the source does, and a struct literal is the
concatenation of its fields. And in VHDL a conditional inside an
expression and a truth value as a bit are functions, \code{mux} and
\code{tobit}, declared in every architecture, since VHDL-2008 has
neither as an operator.

The netlist was checked as the model was: a testbench that could not
fail would prove nothing. The register file is not traced, so the
testbench cannot see it, but every write to it shows in the
writeback output the next tick, and every read shows in the program
counter and the outputs that follow; the six hundred and six
cycles of the demonstration exercise every instruction of
the program through both, the two traps and their returns included,
since a wrong \code{mepc} or \code{mtvec} is a wrong program
counter, and every register of the sequencer at every step of three
multiplies and two divides, since the whole unit is traced. A second
run, \code{vreteno\_amo}, is the same machine with the A extension's
program in place of the demonstration, every AMO, \code{lr.w} and two
\code{sc.w}, and the core's netlist is checked against its trace the
same way, under its own entity's name (issue~1010). The
router, the timer, the serial port and its bridge are checked from
the same trace as entities of their own, each channel's six wires and each
line among the signals, so the bus is checked from every end. One
channel has a name on each side, and the run traces it under one of
them; a port traced under another name than its own says so in the
ports file. The checks that failed on the way, a truth value rendered
ambiguously, a bitwise \code{and} rendered as a comparison, a
concatenation of ambiguous type, a slice of an expression and a
temporary named with a leading underscore, all failed to analyze or
differed at their first tick, which is the test doing its work.

\section{Synthesis}
\label{sec:v-synth}

The Verilog synthesizes, and the build does it: \code{bazel build
//cpu/vreteno:vreteno\_synth} runs Vivado, which \code{rules\_vivado}
installs itself from the AMD archive into a cache the machine's
workspaces share, so the toolchain is as hermetic as nvc and
Verilator are, once per machine rather than per workspace. The target
is an Artix-7, \code{xc7a200tfbg484-2}, with a clock constraint of
ten nanoseconds and no pins, since the question is whether the
netlist synthesizes and to what, not where it goes on a board.
Synthesis took two minutes and finished with no error and no
critical warning. Table~\ref{tab:v-synth} is its report, run on
2026-10-04 at commit \code{a2b06993}.

\begin{table}[htbp]
\caption{The core on an Artix-7, after synthesis; the data memory,
the timer and the serial port, units of their own, are not in it.}
\label{tab:v-synth}
\centering
\begin{tabular}{@{}lr@{}}
\toprule
Resource & Used \\
\midrule
LUTs, as logic & 2624 \\
LUTs, as distributed RAM & 48 \\
Flip-flops & 809 \\
Block RAM tiles & 2 \\
DSP blocks & 4 \\
\bottomrule
\end{tabular}
\end{table}

The memories are where the third stage shows. The instruction
memory, written by nothing and holding the program from its
\code{initial} block, is a read-only memory of 1024 words by 32 bits.
The fetch reads it at two addresses, the word at the counter and the
word after, since an instruction may cross from one into the other,
and the tool makes that two block RAM tiles, one per address, rather
than one tile's two ports; they are the core's two tiles. Without the
compressed instructions the fetch read one word, the memory was one
tile, and the logic was 1952 LUTs: the second read, the choice of
halfword and the expander cost some two hundred LUTs, a tenth of the
core. The data memory, a
device of its own now, is four block RAMs of 1024 by 8 in two tiles
in its unit, each written by a store on its port and read by a load
into the memory's register on the other: the synchronous read a
block RAM has, which is what the third stage was for while the
memory was in the core. With two stages the data memory had been
560 LUTs of
distributed RAM, since the load read it combinationally and the
store read it to merge a byte; the lanes removed the store's read
and the stage registered the load's, and the LUT count fell by a
third. The register file, 32 by 32 with one write port and two
reads, is twelve of the small distributed primitives, which is what
a register file of that size is on this family. The flip-flops are
the core's registers less the bits a mask makes constant: \code{mstatus} has two bits that
can be set, and \code{mtvec} and \code{mepc} lose the alignment
bits the core clears on every write. The rest of the control
registers stay, since every bit can be read back through a CSR
instruction. Synthesis also
trims \code{f\_w1}, the word after the counter as a fetch from the
bus returned it, from 32 bits to 16, with a warning: an instruction
that crosses into that word takes only its lower half.

The netlist is also placed and routed, \code{bazel build
//cpu/vreteno:vreteno\_pnr}, with the same constraint and no pins,
so the bitstream is a stub and the reports are the point. It is
synthesised out of context, so its ports are the partition's pins
and not pads: the core has 308 ports and the package 285 pins, and placement failed on them until it was
(issue~971). Run on 2026-10-04 at commit \code{a2b06993}, the routed
design meets ten nanoseconds with 0.267~ns of slack, and hold with
0.132~ns, in 2622 LUTs and 809 flip-flops. The critical path is
9.717~ns, of which routing is two thirds, through seventeen levels of
logic, nine in dedicated \code{CARRY4} primitives: from the first register field of the
instruction register, through the register file's read and the
operand's forwarding, along the ALU's adder, into the writeback
stage's result. The last run that placed, on an older core, had
9.661~ns, 0.296~ns of slack and hold 0.102~ns, from the second register field
through the operand's forwarding into the ALU's compare and the
writeback value.
The expander of compressed instructions is in the fetch, before the
instruction register, and is not on it. Before the compressed
instructions the slack was 0.315~ns, on a path from the writeback's
destination register, through the forwarding compare into an
operand and the jump's add, to the instruction memory's address.
With the data memory in the core that path had
0.643~ns, and with the step's pieces written in place rather than
as functions 0.776~ns: the synthesis and the placement move by that
much on one design written three ways. Before the device load's wait moved into
writeback the
slack was 0.129~ns, on a path from the same read through the
address add and the decode of whether the address is a device's
into the fetch's hold: the fetch waited on the address. Now the
hold is a register's, and the path is gone. With two stages the
slack had been 0.263~ns, on a path from a register through the
decode and the ALU into the register file's write. With three stages and the
load forwarded it was 0.361~ns, on a load followed at once by a jump
through what it loaded. With the traps, whose return and handler
addresses widened the fetch's multiplexer, that path missed by
0.36~ns. That is what the stall of Section~\ref{sec:v-core} was
for: it took the data memory's read and the load's extension off the
forwarding path, and the slack tripled. The sequencer of the M
extension adds nothing to the path, since its result is a register
and its steps are a cycle each, the multiplier's included, which
runs from registers to registers. The tool had an opinion of its own
on the way. Left to itself it moved three of the four data memory
lanes into distributed RAM to shorten the load's path, so the
board's constraints pin the lanes to block RAM, where the third stage
put them; the data memory is a unit of its own now, and not in the
core synthesised here.

\subsection{The board}

The board is the Alinx AX7A200B, the Artix-7 the part was chosen for.
Everything that computes is one lowered module, \code{board}, the
netlist \code{\#[lower]} writes for \code{vreteno32::board::Board}: a
unit of units holding the core and its tracker, a router of nine, and
the peripherals of Table~\ref{tab:v-map}, each behind a tracker of its
own or, for the page, the interrupt controller and the debug module,
behind an AXI-Lite bridge. The page at \code{0x3000} is seven
peripherals of a sixteenth each, in Table~\ref{tab:v-page}: the
serial port, the pulse width modulator, a slot for a peripheral on a
clock of its own, the remote peripheral whose transactions leave the
board as Ethernet frames, the Ethernet port's registers, the entropy
source, and the configuration flash's SPI master. Both tables are written by \code{//tools/memmap} from
the maps the router and the page's bridge decode, \code{BoardMap} and
\code{SlotMap}, so they are the map the netlist has and not a copy of
it. The interrupt controller's line is the core's external interrupt.
The configuration flash, the chip that holds the bitstream, is reached
twice (issue 312): through that master, for commands such as the
chip's identity, and through a window of 16~MiB at
\code{0x2000\_0000} on the router's eighth port, where an ordinary
read is answered with what the flash holds. Between both and the chip
stand the \code{STARTUPE2} primitive, the only way to the flash's
clock pin once the part is configured, and the pins, which wait for
the end of configuration and refuse write enable, so nothing the core
runs can change the bitstream. On the board, \code{flashid} read the
chip's identity, \code{0020ba18}, the MT25QL128's; found its write
enable latch clear after a write enable the pins refused; and found
the bitstream's sync word through the window at byte 48. Nothing was
written to the flash.
The DDR3 memory is a gigabyte behind AMD's MIG controller, which the
build generates with its own AXI4 port, and \code{AxiPerPins}, which
puts the router's port onto the controller's pins with no tracker in
front, since the controller keeps transactions outstanding itself
(issue~\#1023). The controller is a module from elsewhere, so the
netlist names its wrapper and does not write it, and the memory's pins, the data and strobe lines among
them, are the module's own ports. \code{board\_netlist} writes it with
a program already in the instruction memory and the data memory,
since nothing on the machine loads them at run time.
A program loaded into the DDR3 keeps its stack in 64~KiB at
\code{0x1\_0000} rather than at the top of its region, since the DDR3
is not cached for data and a stack there paid the controller's latency
on every spill and reload (issue~\#1278). That window was at first a
memory on the router's ninth port, the data memory's unit made larger:
a frame of the GL icosahedron's list in \code{board\_test} took 5.42
million cycles with its stack in the DDR3 and 3.30 million with it
there. A load from it still crossed the router and a tracker both ways
and took 9.14 cycles, about 30\% of the frame.

So the window is now a data RAM in the core, on its own port and not
on the bus (issue~\#1275), four lanes of a byte, written at one
address and read at another, as a block RAM is. Whether an access is
in the window is found from the upper half of its base register and
the carry out of the lower halves' sum, since the offset is twelve
bits, signed; that sum also gives the lanes their address, so the
decision costs the bus's valid half the adder's carry chain. A load
there does not go to the bus: the lanes read at the address the access
had in execute, into a register, and writeback takes the word from it
in the next cycle, so a stack load costs what any load costs when the
next instruction does not use it. A store writes the lanes in the
cycle it runs, under its strobes; \code{lr.w}, \code{sc.w} and the
AMOs work there as on the bus. The GL frame takes 1.54 million cycles,
31\% fewer than the 2.23 million it took with the window on the
router. Nothing but the core reaches the window now: not the debug
module's access to the system bus, so a debugger reads a program's
stack through the core (OpenOCD's \code{riscv set\_mem\_access progbuf
sysbus}), and neither Razboj nor a DMA engine. Nothing is fetched from
it either; a fetch there is refused, in the model as on the board.
The netlist the documents simulate and the layout maps keeps four
words a lane, where the window's addresses wrap, since a cell library
has no RAM and 64~KiB of flip-flops would be most of the layout; the
board's is 16384.

The data memory's tracker hands a request to the memory, and a read
beat to the router, over unregistered channels (issue~\#1291), so each
crosses in the cycle it is offered: a load from the data memory takes
two cycles fewer. The tracker runs before the memory and the router in
the board's join, as the sender of an unregistered channel has to.

\begin{table}[htbp]
\caption{The board's address map, lowest first, as the router decodes
it.}
\label{tab:v-map}
\mmget{board}{vreteno}
\end{table}

\begin{table}[htbp]
\caption{The peripheral page at \code{0x3000}, as its bridge decodes
it.}
\label{tab:v-page}
\mmget{page}{vreteno}
\end{table}

The peripherals that declare their registers with \code{regmap!} are
listed by the same tool, each from the base the maps above give it:
the timer in Table~\ref{tab:v-timer}, the serial port in
Table~\ref{tab:v-uart}, the pulse width modulator in
Table~\ref{tab:v-pwm}, the Ethernet port's registers in
Table~\ref{tab:v-ethslots}, the entropy source in
Table~\ref{tab:v-trng}, and the configuration flash's master in
Table~\ref{tab:v-spi}. That is every peripheral on the page with
registers of its own: the remote peripheral's behaviour is a program on
another machine, and the third slot is brought out of the unit, where
the flagship puts its video peripheral, whose table is in the flagship
document.

\begin{table}[htbp]
\caption{The timer's registers.}
\label{tab:v-timer}
\mmget{regs}{timer}
\end{table}

\begin{table}[htbp]
\caption{The serial port's registers.}
\label{tab:v-uart}
\mmget{regs}{uart}
\end{table}

\begin{table}[htbp]
\caption{The pulse width modulator's registers.}
\label{tab:v-pwm}
\mmget{regs}{pwm}
\end{table}

\begin{table}[htbp]
\caption{The Ethernet port's registers.}
\label{tab:v-ethslots}
\mmget{regs}{ethslots}
\end{table}

\begin{table}[htbp]
\caption{The entropy source's registers.}
\label{tab:v-trng}
\mmget{regs}{trng}
\end{table}

\begin{table}[htbp]
\caption{The configuration flash's SPI master's registers.}
\label{tab:v-spi}
\mmget{regs}{spi}
\end{table}

Listing~\ref{lst:v-board} is the one file written by hand, and it
holds only what a board needs and a lowered unit cannot say. The
controller makes the clocks: it takes the board's 200~MHz, which is
also its delay reference, runs the memory at 400~MHz, and hands back
the 100~MHz the design runs on, with a reset that holds until that
clock is good, so the board has no clock generator of its own. The
controller's reset is a pulse at power-on and nothing else; the core's
follows the button, the controller's clock being good and its
calibration, so the core starts only once the memory it may reach is
ready. The LEDs, lit when driven low, show the
core halted, the serial line driven at least once, the memory
calibrated, and a heartbeat. The second LED held the heartbeat until
the board was first run, and the heartbeat moved to the fourth, which
held the clock generator's lock: a blinking heartbeat says as much,
since the counter behind it runs on the clock the controller makes. What no LED said was
whether a byte ever left the core's serial port. That question is
issue~195, where the board halts with every LED on and the host's
serial port stays silent, so the second LED now latches the first start
bit on the line. A heartbeat in a new place also tells somebody at the
board that a new bitstream is in the part.

The program for the board says more than the one a simulation runs.
Built with \code{--cfg=board\_run} it puts a header line out before
it touches the memory, and a dot every ten seconds for half a minute
before the test and two minutes after it, so that a reader who
attaches at any moment can tell a silent serial path from a core that
never ran. It also halts only when every word came back, and spins
when one did not, which puts the verdict on the first LED for anybody
who cannot read the line. Simulation runs the program without those,
since tests directly verify the entire output stream and every byte is
simulated time they pay for.

The integrated design is simulated under Vivado's simulator by
\code{//cpu/vreteno/board/sim:board\_test}: the top, the netlist with
the DDR3 test in its memories, the controller, and two Micron models
of the board's chips on its pins. The program writes sixteen words
across the memory's region and reads them back. In the run the memory
calibrates at 108.1~microseconds, the controller's fast calibration
being what its simulation variant provides, and the core says
\code{ddr3 ok} on the serial port and halts at 130.4, which is what
the test's terminal reads. The controller alone is checked the same
way, against the same models, by \code{//ddr3/sim:wrapper\_test}, and
AMD's own example design, its traffic generator writing and reading
the memory through the controller, by \code{//ddr3/sim:example\_test}.

\code{bazel build //cpu/vreteno:vreteno\_board\_pnr} synthesises it,
places and routes it with the board's pins and the memory's, and
writes the bitstream. Every timing constraint is
met: the worst setup slack is 0.307~ns, the worst hold slack
0.033~ns, over 32\,815 endpoints on the clocks the controller's
generators derive from the one constraint on the board's 200~MHz
input. No register is without a clock, and no internal endpoint is
unconstrained. The design takes 12\,296 LUTs, 11\,209 flip-flops, five
block RAM tiles, four DSP blocks, 86 pins, the controller's PLL and
MMCM, its 64 output and 32 input serialisers with 32 delays, and the
\code{STARTUPE2} that clocks the configuration flash. These are the
numbers of \code{main} at eed5b7f, with the configuration flash and
the entropy source's capture buffer on it (issue 842); the pins'
constraints, the MDIO master, the debug transport and the SD card
host came after, and the size has not been measured since. The
memory's pins specify no input or output delay, as in the a200t
examples, since the controller calibrates their timing itself. The
board's other pins are timed, or cut by a false path, in
\code{board/ax7a200\_io.xdc} (issue 847), so none is left without a
delay. The report's
eight combinational loops are the entropy source's eight rings, which
are loops by design and cut by a false path.

The run on the board below was made with the design before the DDR3
memory was on the bus, when the top wired the units to each other by
hand around an MMCM. The design with the memory has run since, first
with UberDDR3 and then with AMD's controller; Section~\ref{sec:v-next}
says when and what it said.

The board sits on another machine, with its programming cable and
its serial bridge, and is programmed from here. Vivado's
\code{hw\_server} and the cable's libraries come out of the
hermetic installation, \code{vivado\_extract}, into a bundle that
\code{bazel run //cpu/vreteno/board/remote:hw\_server} uploads over
ssh, starts there with an idle timeout, and tunnels back to a local
port; \code{bazel run //cpu/vreteno:vreteno\_board\_prog} then has
Vivado connect to that port and write the FPGA, which the next power
cycle forgets, and \code{vreteno\_board\_flash} the same way writes
the QSPI flash, which it does not. The serial port is watched the
same way: \code{bazel run //cpu/vreteno/board/remote:serial} uploads
a small program, one static binary in Go on the standard library
alone, and runs it on the far machine over ssh and under a timeout
that frees the port whatever happens here; it prints what the board
says for a while and may type a reply once the board has said a
line.

The watcher is opened before the programming, since the board's
program says its first line within seconds of the part being
configured. On 2026-10-04 a watcher opened before
\code{vreteno\_board\_prog} and held across the reconfiguration
received the DDR3 test whole (issue~965). It was not always so. Issue~195
measured the opposite: with a bitstream that says one line a second in
the part, a watcher held open across a reconfiguration saw no bytes in
150 seconds, and one opened straight after saw 360 bytes in 30, which
was read as the reconfiguration's high-impedance pins making a break on
the board's serial bridge. That cost an evening of runs that looked
like a silent design, which is why the board's program says a header
and a dot every ten seconds: a board that keeps talking for two minutes
can be read whenever the reader is ready, whichever order turns out to
hold.

The bitstream in that measurement is \code{//cpu/vreteno/board/probe},
which has no core, no memory and no bus: the board's 200~MHz clock, a
counter, and a transmitter that says one line a second on the same pins
and the same I/O standards the design uses. It is the control for any
question about the serial port, since a line from it and silence from
the design says the fault is in the design, and silence from both says
it is the path or the reader.

The board did what the runs did. Programmed over the tunnel, it said
OK and a newline on the serial port at 115\,200 baud, the watcher
typed \code{yes}, and the board echoed the three bytes: six bytes on
the line, the same six the lockstep test's model had recorded. The
LEDs were looked at on the board itself, by an author with a camera:
they did what that run's top said, and the photographs and the video
are the record. The top below is today's, whose LEDs show the halt,
the first start bit on the serial line, the memory's calibration and
a heartbeat. Holding the reset puts the halt LED out; releasing it says OK
and waits for three keys again.

\lstinputlisting[caption={\code{board/vreteno\_board.v}: the board
top, by hand.},label={lst:v-board}]{cpu/vreteno/board/vreteno_board.v}


\subsection{The board with a logic analyser}
\label{sec:v-ila}

A design on the board shows four LEDs and a serial line, and a
question about a wire inside it had no answer but another run under
Vivado's simulator. The analysis in \code{docs/on-chip-debug.md}
lists the ways to see more and puts an integrated logic analyser
first, because the rules exist and the design does not change.
\code{//cpu/vreteno:vreteno\_board\_ila\_pnr} is that bitstream
(issue~239): the same lowered module under a top that is
\code{vreteno\_board.v} with one block added, Vivado's analyser
generated by \code{vivado\_ila} with thirteen probes and 4096 samples,
sampled on the core's clock. It is built beside the shipping
bitstream and is not it, since the analyser costs block RAM and a
run of its own: routed, the design takes 11\,010 look-up tables,
11\,334 registers and 19 block RAM tiles, against the 8\,100, 6\,959
and four of the board it was measured beside, the one with UberDDR3, and meets its timing with 0.6~ns to spare.

The probes are the nets a reader of the board asks about first, and
the top says why each is there: the core's reset, its halt and the
memory's calibration; the serial line both ways and the break that
resets the core; the modulator's four channels; the third slot's
three request channels, on which a program that touches
\code{0x3200} on this board is seen waiting for an answer that never
comes; the remote peripheral's frames, byte by byte; and the two
flags the LEDs show. The core's retire stream is not among them: it
ends inside the lowered module and is not a port of it, which is what
the retire buffer of issue~240 is for. The memory's command pins are
on the 400~MHz clock and are left out rather than sampled across a
clock boundary.

A capture is read back over the JTAG cable through the probe file
the implementation writes beside the bitstream, with
\code{//cpu/vreteno:vreteno\_board\_ila\_read}, once the board has
been programmed with \code{vreteno\_board\_ila\_prog}. That was done
on September 22, 2026, in a board window with nobody at the board:
the analyser bitstream was programmed over the tunnel, the core was
run with no trigger, so it captured the board as it stood, and the
4096 samples came back as the VCD in Listing~\ref{lst:v-ila}, kept
in the tree as \code{docs/ila\_capture.vcd}. Nothing in it changes
for 41~$\mu$s, which is the board at rest a minute after
configuration, and every probe says what the LEDs say and more:
\code{calib} high, the memory calibrated; \code{halt} low and
\code{said} high, the core running and the serial line driven at
least once, which is the program having said its line and waiting
for a key; \code{rst} and \code{brk} low, \code{uart\_tx} and
\code{rx\_sync} idle high; the third slot's three channels with no
valid, holding the last phases that crossed them; the remote
peripheral's frame channel idle; the modulator's four channels low.
The analyser's own three marks, \code{\_TRIGGER}, \code{\_WINDOW}
and \code{\_GAP}, are Vivado's, and the trigger mark at sample 2048
is where the run began.

The read went through Vivado's Tcl by hand rather than through the
rule's generated script, which failed twice before it reached the
part: it looked for its log library under a name that exists in no
repository, and it asked for \code{arm\_hw\_ila}, which is not a
Vivado command. Both are fixed upstream in
\code{bazel\_rules\_vivado}, with the Tcl the rule generates and the
Tcl run by hand being the same after the fix.

\lstinputlisting[style=ascii,caption={\code{docs/ila\_capture.vcd}: the capture,
as Vivado wrote it.},label={lst:v-ila}]{ila_capture.vcd}

\subsection{The bus from the cable}
\label{sec:v-jtag}

The loader takes a program over the serial port, and the analyser
shows a capture, but neither reaches the bus from outside: a program
that has hung, or a peripheral that is not answering, could be read
only by the core, which is the thing in question. The analysis in
\code{docs/on-chip-debug.md} puts a host path onto the bus third,
after the analyser and before a debug module, and this is that path
(issue~\#241): Vivado's JTAG-to-AXI master, generated by
\code{vivado\_ip} as \code{//cpu/vreteno:jtag\_axi}, on a second host
port of the lowered board, reached over the cable that programs the
part.

The board changed in two places for it. Its ports are two structs
now, \code{BoardIn} and \code{BoardOut}, the old ports under their old
names and the master's pins under \code{jtag\_}, one field of type
\code{AxiHostPins}, passed whole to an \code{AxiPins} that joins them
to the link's channels. And the router got an arbiter in
front of it, \code{Arbiter2} then, with the core's tracker on one port
and the pins on the other. An arbiter puts the port number above a
host's own identifier, so the peripheral side of the bus allocates four
bits where it previously allocated two: two for a host's own identifier and two for the port,
which is room for four hosts. Four filled it, \code{Arbiter4} over
the core, the pins and the Ethernet port's two engines
(issue~\#151); the fifth widened it, below. The debug
transport of Section~\ref{sec:v-dtm} does not: it shares the pins'
port, through an \code{Arbiter2} of its own in front of it, the pins
and the transport one bit of identifier each, which is what the
JTAG-to-AXI master's IP ever used. Every peripheral's
tracker, the memories, the timer and the bridges were widened with
it; none of them names the width except at its instantiation.

The video scanout's fetch is a fifth host (issue~\#151). It first
shared the Ethernet send engine's port through a nested
\code{Arbiter2}, as the debug transport shares the pins', so that
nothing past the arbiter widened. With more hosts coming, the SD
card's engine and Razboj among them, the choice was to widen once and
stop nesting (issue~\#1022): the peripheral side allocates five bits of
identifier, two for a host's own and three for the port. Three bits of
port number cover up to eight hosts, so each one after the fifth
raises the arbiter's count and widens nothing behind it. Every
peripheral's tracker, the router, the bridges and the DDR3's were
widened again with it, each at its instantiation.

The SD card's engines took the sixth port (issue~\#912), and the
arbiter is \code{Arbiter6}. They are two, a store engine for a read
and a fetch engine for a write, and they reach that port through an
\code{Arbiter2} of their own. That is a nest after all, and the reason
against nesting does not apply to it: hosts behind one port share its
turns, and these two never offer at the same time, since a command
is a read or a write and none starts until the last has finished.
The host checks the two busy lines are never high together. The nest
keeps two ports free, for Razboj and a texture fetch.

Razboj took the seventh (issue~\#985), and the arbiter is
\code{Arbiter7}. Its rasteriser draws into the frame the scanout
shows, at \code{0x4200\_0000}, from a display list at
\code{0x4280\_0000}. It reads the list's count from a doorbell on the
peripheral page's tenth slot, at \code{0x3900}: a register that holds
zero from reset, so no word of memory can start a list nobody wrote.
The doorbell also drives a line to the rasteriser, high while the
count is not zero, and the rasteriser reads the count only then. An
idle rasteriser therefore puts nothing on the bus. Reading the count
back to back, or even once in 256 cycles, moved the timing of every
other host's accesses, and the board tests that time the DDR3 path
or sample the entropy source failed with it. Razboj's rows are
$2^{10}$ words, so the scanout's stride became 4096 bytes, of which
it shows the first 640 words. A board test writes a list of one rectangle,
rings the doorbell, and reads the pixels back over the bus.
On the board, \code{razprobe} does the same with a full screen and
three shapes, shows them through the scanout, and says on the serial
port whether the pixels it read back are the list's
(issue~\#1169).

The cost of a host is its share of the bus when every host is busy.
For the scanout, a line of 640 words, 2560 bytes, must arrive within
the 31.8\,$\mu$s of the line before it, about 80\,MB/s. The link
moves at most four bytes a beat at 100\,MHz, 400\,MB/s, and the DDR3
controller behind it less. With every port offering a burst, each
wins one turn in seven, and a turn is a burst of sixteen beats. The
SD port offers a burst only as fast as the card fills one, though,
at most 12.5\,MB/s on four lines at 25\,MHz, so the other six share
nearly all of the link and the scanout's worst share is about
65\,MB/s, short of a line's need. That bound holds only while
every other host moves bursts back to back for a whole line, which
the core's single-word accesses do not. The tighter limit is behind
the arbiter: the DDR3 path takes one access at a time, and measured
at the model's best latency the whole path moves about 79\,MB/s
(issue~\#1023), so a scanout under load already fills it. More than
one access in flight, or whole bursts to the controller, is what
changes that, and is issue~\#1023's to decide.
The scanout's sticky underflow bit is what says whether it ever
happens, and the board check runs the scanout with the core copying
in the DDR3 and the Ethernet port sending, so that the bit is read
under contention and not on an idle bus.

On the board it did happen (issue~\#1209): under that load the
picture showed short streaks of colour from elsewhere on the row, the
line buffer running dry and showing what the line before had left. A
burst of sixteen paid the path's latency and a turn at the arbiter
forty times a line, and the core's single words between them are turns
too. In the board's simulation, a line as long as the board's, 3175
cycles, took 1964 idle and 4682 under the load, every line late. The
scanout now fetches a line in ten bursts of 64 beats: a line from a
4~KiB-aligned base never crosses a 4~KiB boundary, and
\code{LineFetch} cuts a burst at one if it would. The same runs then
take 974 idle, 1206 under the load and 1200 with Razboj drawing as
well, under two fifths of a line, and
\code{the\_scanout\_keeps\_up\_with\_the\_boards\_load} holds every
line to half of one. Longer bursts fetch faster, 942 under the load
at 128 beats, but a burst holds the path, and a load of the core's
behind one waits for it: 88 cycles with no scanout, 94 beside bursts
of 64 and 153 beside bursts of 128, which
\code{a\_core\_load\_waits\_behind\_the\_scanout} measures. A column whose word is late anyway now shows
magenta, \code{LATE}, rather than a stale word, so an underflow looks
like one. And a late line's words land in its own row: before, the
words of a line that came after its start went into the half filling
for the next row, which then showed them as its own, one row low
(issue~\#1233).

\code{vreteno\_board\_jtag.v} is the top with the master, the board's
own top with one block added, and every other top ties the pins off,
so the shipping bitstreams are unchanged in what they do.
\code{//cpu/vreteno:vreteno\_board\_jtag\_probe} is the proof
(Listing~\ref{lst:v-jtag}): from the hardware manager, words go into
the data memory and the DDR3 and come back, and the timer's
\code{mtime} is read twice and differs, which is a peripheral
answering while the core runs its own program.

\lstinputlisting[style=ascii,caption={\code{board/jtag/probe.tcl}:
the bus from the cable, and what it checks.},label={lst:v-jtag}]{cpu/vreteno/board/jtag/probe.tcl}

It ran on September 22, 2026, in a board window with nobody at the
board, and every word came back (Listing~\ref{lst:v-jtag-run}, kept
in the tree as \code{docs/jtag\_probe.txt}): two words into the data
memory and two into the DDR3, read back as written, and \code{mtime}
two different numbers a moment apart. The core was running the
board's own program from the boot memory the whole time; the master
and the core share the bus through the arbiter, and neither knows
the other is there. Routed, the design takes 10\,795 look-up tables and 11\,244
registers, against the 8\,100 and 6\,959 of the board it was measured
beside, the one with UberDDR3, and meets its
timing with 0.17~ns to spare, closer than the board alone, since the
master and the arbiter add a stage on the bus's paths.

\lstinputlisting[style=ascii,caption={\code{docs/jtag\_probe.txt}: what
the probe said.},label={lst:v-jtag-run}]{jtag_probe.txt}
\subsection{The debug module}
\label{sec:v-dm}

The core's debug mode, Section~\ref{sec:v-core}, is reached from
outside through the debug module, the RISC-V debug specification's
module as far as a debugger on the bus needs it: the second of the
three steps of issue 154, after the mode and before a transport of
its own. The specification puts the module behind a small address
space of registers a transport reaches from the cable; here that
space is on the bus, an AXI-Lite peripheral behind a bridge of its
own on the router's seventh port, the 64~KiB from \code{0x1000\_0000},
each register a word at four times its number. So the JTAG host of
the last section reaches it from the hardware manager today, and a
transport reaches it the same way when there is one. \code{dmcontrol}
encodes the halt and resume requests and \code{dmactive};
\code{dmstatus} says whether the core is halted or running and
whether a resume was acknowledged; \code{abstractcs}, \code{command}
and \code{data0} are the abstract command, access-register only, one
32-bit word: a general register, which the core reads and writes for
the module through its register file's lent ports, or a CSR, which
it reads through its CSR read and writes only for \code{dpc} and
\code{dcsr}. A command while the core runs is \code{cmderr}~4, an
unsupported one 2, one while another is in flight 1, and a command
under a standing error is ignored, as the specification says. There
is no program buffer. The specification's system bus access,
\code{sbcs}, \code{sbaddress0} and \code{sbdata0}, is served beside
the module by the transport's bridge, where the bus master is
(Section~\ref{sec:v-dtm}); a debugger on the board's own bus, as the
JTAG host is, needs none of it.

The module's own test plays the core, sets halted and the word read,
and reads the lines back; the board test puts a debugger on the JTAG
pins, a master model that follows a plan of writes, reads and waits,
against a program counting to three thousand: it halts the core, reads
the count and \code{dpc}, writes the count to its end, resumes, and the
program finishes at once, which is the write having landed. On the
board, \code{vreteno\_board\_dm\_probe} is to do the same from the hardware
manager against whatever program runs; that run is owed (issue 154).

\subsection{The debug transport}
\label{sec:v-dtm}

The third step of issue 154 is the transport, so that a stock
OpenOCD, and gdb through it, reaches the module from the cable that
programs the part, with nothing written for either of them.
\code{Dtm}, in \code{txhdl\_parts::dtm}, is the specification's
transport, version 0.13. It sits behind \code{BSCANE2} on the
device's fourth user scan chain, which every top instantiates as
\code{bscan\_user4}, and it runs on the cable's clock. OpenOCD
reaches it through its BSCAN tunnel, \code{riscv use\_bscan\_tunnel
5}, which encapsulates each of the transport's own scans inside one data
scan of that chain. The examples document has the transport alone
and the tunnel's framing.

Each access the transport takes crosses to the board's clock through
a \code{ChanCdc}. There \code{DtmBridge} makes it a read or a write on
the link, through a host tracker of its own: one of the module's
registers at \code{0x1000\_0000} plus four times its number, or the
system bus access, which the bridge serves itself. It serves the access
8, 16 or 32 bits wide, a narrow one as a word on the link with only
its lanes' strobes set, since gdb reads and writes a compressed
instruction, and the \code{c.ebreak} it puts over one, a halfword at a
time (issue~873). The bridge's host
and the JTAG master's pins share the arbiter's second port through
an \code{Arbiter2}, taking turns, so nothing past the arbiter changed.

Two tests check it on the whole board in simulation, with a model of
the Artix-7's TAP in front of \code{BSCANE2}, its instruction
register six bits and \code{USER4} at \code{0x23}. In the first,
OpenOCD~0.12.0 itself, pinned and unmodified, drives the board over
its \code{remote\_bitbang} adapter. It finds the device, examines the
hart (one hart, 32 bits, \code{misa} \code{0x40001104}), halts it,
reads \code{pc}, reads four words of the data memory by system bus
access, steps one instruction, writes a word and reads it back, and
resumes it. The run takes under half a minute, so it is
in the suite. Stepping is what found the module refusing a write of
\code{tselect}: OpenOCD counts the triggers that way, and refused, it
stops writing any CSR by the abstract command, \code{dcsr} and its
step among them. With no trigger module the module now takes that
write and changes nothing. In the second test, the JTAG master reads the data memory
without a pause while the transport reads it too, and both get their
words: the master read 122 words, 79 of them while the transport's
access was in flight, so neither port waits for the other to stop.

\section{Programs in Rust and C++}
\label{sec:v-rust}

Everything above runs programs written with the assembler in
\code{program.rs}, which is enough to check a core and no way to
write software. The core is an RV32IMAC machine, so a compiler can
target it, and the toolchain for that is fetched by the build like
every other: \code{MODULE.bazel} names the bare-metal RISC-V target
beside the host's, and the ruleset downloads that target's standard
library by checksum. The linker is \code{rust-lld}, which comes with
the toolchain, so there is no cross binutils to install either.
Nothing is on the machine that Bazel did not put there.

A program for the core is built for another machine than the one the
build runs on, so \code{vreteno\_image} bridges the transition: a rule
with a platform transition that builds the binary for the core, runs
\code{//tools/elf2vreteno} over the linked ELF, and gives back the
image as Rust source. The rest of the tree depends on that source, so
a program for the core is built by \code{bazel build //...} like
anything else, with no flag to remember.

The transition also sets the flags every crate built for the core is
compiled with: optimised for size, no debug assertions, no overflow
checks. A library a program links is built with them as the program
is, and a build of the same library for this machine, its tests
included, keeps its own (issue 1248). Before, a library that named no
flags of its own was built for the core unoptimised. The logo's pixel
table was copied whole for each pixel read, and the paint took 31
seconds where it takes two (issue 1245). The GL ES library's list of
a frame of the icosahedron took four times the instructions it takes
now, and the program twice the bytes.

The board's own test, \code{//cpu/vreteno:board\_test}, is built the
other way round: everything it links, the simulation of the whole
board among it, is optimised at level 3, through \code{opt\_test} in
\code{tools/opt\_test.bzl}, with its debug assertions and overflow
checks kept on, so it checks what it checked at level 0 (issue 1250).
At level 0 the board ran some 1500 cycles a second and the test took
670 to 1100 seconds; optimised it takes about 35.

\subsection*{What the machine will not let a compiler do}

Two of the core's properties decided how a program had to be
written, and each of them the core met by going quietly wrong rather
than by refusing. One of them is gone.

The first was the one that shaped the linker script. A load below the
data base never left the core: no request was sent, and the
instruction that was waiting retired with whatever word the previous
device load had left in its register. The instruction memory was not
a device on the bus, so a constant placed there was a constant the
program could not read, and it was not an error, it was a wrong
answer. So \code{.rodata} was linked into the data memory rather than
beside the code. Since issue 268 the boot memory is on the bus as
well, read-only at address zero (Section~\ref{sec:v-fetch}), every
load goes to the bus, and \code{.rodata} is linked beside the code
where the linker script puts it; the data memory stores only what a
program writes, and \code{Dmem::with} gives it its initial bytes.

The second is the size. A program is 4096 bytes, 1024 words, the
loader truncates a longer one without saying so, and the fetch takes
its word by \code{pc[11:2]}, so a program that runs off the end wraps
into itself. Optimising for size is not a preference here. It is also why
debug assertions are off: one of them pulls in \code{core}'s panic
formatting, which is larger than the whole instruction memory. The
compressed instructions help: the target the ruleset maps has them,
the core runs them, and the greeting is 25 words where it was 33
without them.

Two more things nothing else does: the stack pointer is zero at
reset, and \code{.bss} is not cleared. An entry stub does both, in
the assembly that has to come before Rust can run at all. And a
program ends by writing an odd value to \code{mhalt},
\code{csrwi 0x7c0, 1}, which is one instruction and needs no runtime;
\code{ebreak} is a breakpoint here as the specification says, and
traps. Neither is in the greeting any more: they are
\code{vreteno\_hal}'s, Section~\ref{sec:v-hal}, and the greeting is
the program and nothing else.

\lstinputlisting[caption={\code{cpu/vreteno/rust/hello.rs}: Rust
compiled for the core. Run the whole machine on it with
\code{bazel run //cpu/vreteno:hello}.},label={lst:v-hello}]{cpu/vreteno/rust/hello.rs}

\lstinputlisting[caption={\code{cpu/vreteno/rust/vreteno.ld}: where
each section goes, and why it has to go
there.},label={lst:v-ld}]{cpu/vreteno/rust/vreteno.ld}

\subsection*{What every program would otherwise write again}
\label{sec:v-hal}

The greeting was the right amount of machinery for a greeting and the
wrong place to start the fourth program from: a naked \code{\_start},
addresses copied from the board's map, a panic handler, and, in the
program that takes input by interrupt, a trap vector that saved and
restored sixteen registers by hand. \code{vreteno\_hal}, the crate
under \code{cpu/vreteno/rust/hal}, holds that once.

Listing~\ref{lst:v-hal} is the crate less its trap entry: the address
map in one place; the serial port, the timer, the interrupt
controller, the pulse width modulator and the video peripheral as
types whose methods are volatile accesses at the map's addresses and
nothing else; the control and status registers a program touches;
\code{entry!}, which owns the reset path; and \code{halt}. The crate
is inline RISC-V assembly from end to end, so it builds only for the
core, through the transition a program's image applies, and a program
names it in \code{deps}.

\lstinputlisting[caption={\code{cpu/vreteno/rust/hal/lib.rs}: the
map, the peripherals as types, the registers, the entry and the
halt.},label={lst:v-hal}]{cpu/vreteno/rust/hal/lib.rs}

Listing~\ref{lst:v-trap} is the trap entry and the dispatcher. The
entry saves every register a call may change and \code{mepc} into a
frame on the stack, hands the frame to the dispatcher, and on the way
back restores the registers and \code{mepc} from the frame and returns
with \code{mret}. So a handler is an ordinary Rust function that takes
the frame, and one that wants to resume after a faulting or calling
instruction moves \code{frame.mepc} on. The dispatcher reads
\code{mcause} and calls the handler the program registered for that
interrupt line or that exception cause; an interrupt nobody registered
for is taken and ignored, and an exception nobody registered for halts
the core, since a program that traps where it did not expect to has
nothing sensible to resume. Nothing is cleared after the handler: the
core's pending bit for the external interrupt is the controller's
line, as the specification has it.

\lstinputlisting[caption={\code{cpu/vreteno/rust/hal/trap.rs}: the
frame, the entry, the dispatcher, and how a handler is
registered.},label={lst:v-trap}]{cpu/vreteno/rust/hal/trap.rs}

Listing~\ref{lst:v-traps} exercises it. The program registers a
handler for \code{ecall} that counts the call and moves \code{mepc}
past it, makes three calls, registers a handler for the timer's line,
asks the timer for an interrupt four hundred cycles on, turns
interrupts on, and waits in a loop that touches nothing until the
handler has run and disarmed the timer. It says
\code{ecall 3 timer 1}, and \code{//cpu/vreteno:traps\_test} runs it
on the machine of Section~\ref{sec:v-rust} and checks that it did.
The greeting and \code{irq.rs} run on the crate as well, and their
tests are unchanged, which is what says the crate is the same
machinery and not different machinery.

\lstinputlisting[caption={\code{cpu/vreteno/rust/traps.rs}: an
environment call resumed past, and a timer interrupt taken, through
the dispatcher.},label={lst:v-traps}]{cpu/vreteno/rust/traps.rs}

The run is the check. \code{//cpu/vreteno:hello} builds the image,
puts it and its constants in the boot memory and its initialised data
in the data memory, runs the whole machine with the router, the timer and
the serial port, and reads the port's line with the same terminal the
demonstration uses. It asserts what came out, so a program that says
nothing fails rather than passes quietly, and it asserts that the
core halted itself. The same run is \code{//cpu/vreteno:hello\_test},
so the toolchain and the core are checked together by
\code{bazel test //...}.

\subsection*{And in C++}

The other toolchain is the GNU one for bare-metal RISC-V, fetched by
checksum like the first. It is the GNU one rather than a Clang
because of what comes with it: its \code{rv32imc/ilp32} multilib,
RV32IMC and the soft-float ABI, is this core less the atomic
instructions, so its code runs here directly, and it includes
\code{newlib} and a precompiled \code{libstdc++},
so a standard library is something the
toolchain provides rather than something the build would have to compile.

Everything the machine forbids forbids it just as firmly in C++, and
the same linker script says so, since what a machine allows does not
depend on the language above it. Four flags are the language's own.
There is no unwinder and nothing to unwind to, so exceptions are off;
nothing reads a type at run time, so \code{RTTI} is off; a static
with a guard would call a lock that does not exist, so thread-safe
statics are off; and the entry point is the core's reset vector
rather than a C runtime's, so the toolchain's startup files are left
out and the reset stub sets the stack pointer and clears the
uninitialised data itself, as the Rust one does.

Deploying a standard library on a machine with 4~KiB of
instruction memory imposes practical boundaries. The portion
consisting of templates and constants is usable, and the program uses
it: \code{std::array}, \code{std::string\_view}, \code{std::for\_each},
and a sum of squares the compiler works out before the core is
started. The half that needs an operating system underneath, a heap,
file descriptors, an unwinder, is not usable, and no toolchain fetch
makes it so. That is freestanding C++: most of the language and a
good deal of the library, and none of the runtime.

\lstinputlisting[caption={\code{cpu/vreteno/cpp/hello.cc}: the same
greeting, in C++, against the standard library the toolchain provides
for this core.},label={lst:v-hellocc}]{cpu/vreteno/cpp/hello.cc}

The run is checked the same way, by \code{//cpu/vreteno:hello\_cc}
and its test. Both languages end at the same place: the machine of
\code{run.rs}, the same terminal on the serial line, and the same
two assertions. What changes between them is the language and
nothing else.

\section{Two of It, Compared}
\label{sec:v-pair}

A machine that must not be quietly wrong runs two of itself and
compares. \code{Pair} is that: two \code{Vreteno} cores on one bus,
with everything they say compared and one line, \code{differs}, that
goes up when they stop agreeing and stays up until reset.

What makes the comparison worth anything is that both cores see
exactly the same inputs in the same cycles. The wires are given to
both; each of the three channels the pair receives goes through a
\code{Tee}, a part that offers each word to both cores and takes the next
only once both have taken it, so neither core can be more than a
word ahead. Each of the three
channels the cores drive goes through a \code{Check}, which takes a
word from both in the same cycle, passes the first core's on and
raises its line when the two differ. Both parts are in
\code{//lib/parts} under \code{redundant}.

The three wires, \code{halt}, \code{instr} and the writeback, have no
handshake, so \code{Watch} compares them every cycle and passes the
first core's on; it also gathers the three checks' lines into
\code{differs}, since a unit of units has no logic of its own to do it
with.

What leaves the pair is the first core's, so a design puts a pair
where it would have put a core and wires nothing differently. What it
adds is one input and one output: \code{fault}, which holds the second
core in reset while it is high, and \code{differs}.

\code{fault} is the self-test. A comparator nobody has ever seen fire
is a comparator nobody knows the state of, and holding one core in
reset for a cycle is a divergence with a known cause that runs through
every part of the comparison path. \code{//cpu/vreteno:pair\_test}
does both halves: eight random programs run with the two cores agreeing
throughout, and then one of them run again with a single cycle of
\code{fault}, after which the line is up. A third run reads the line
again at the end of six hundred cycles and finds it still up, which is
the half of the promise the first two do not check.

What the pair does not do is decide what to do about a disagreement.
It says that there was one; a design around it chooses between a
reset, a trap and a stop, and that choice belongs to the system rather
than to the pair.

\section{What Comes Next}
\label{sec:v-next}

More clock. What sets the timing now is stated in
Section~\ref{sec:v-synth}: the branch on a register, three quarters
routing, which a fourth stage, or a register file read a cycle
earlier, would shorten, and which the bus took the timer's store off
by taking the timer out of the core; then the multiplier, which a
second register in the middle of its adder pipeline, one more step of
the sequencer, would. The lockstep test is the check for either,
since it compares architectural state when an instruction retires.

The design with the DDR3 memory ran on the board on September 18,
2026, and said \code{ddr3 ok}: sixteen words written across the
memory's range and read back, every one of them correct. It ran twice,
once written into the part over JTAG and once from the board's own QSPI
flash after a power cycle, with nothing attached but power. So the
whole path is proven on hardware: the core, the AXI link, \code{AxiWb},
the controller and the board's memory chips. That run was made with
UberDDR3 as the controller; the board integrates AMD's MIG since issue
188, GPL-3.0 having no place in a design this tree claims as
Apache-2.0. The design with MIG ran on September 28, 2026, written into
the part over JTAG, and said \code{ddr3 ok} the same way: the
controller calibrates, the memory returns what it was given, and the
core runs on after. The run from the flash after a power cycle is still
owed (issue~\#188). Since issue~\#1023 the core reaches the controller
through the controller's own AXI4 port rather than through
\code{AxiWb}, and the run of that path on the board is the issue's
fourth step.

\begin{thebibliography}{10}

\bibitem{embedding}
F.~Filmar and D.~Janković, ``Deleting a language: Embedding TxHDL in
Rust,'' project repository \code{HDL/txhdl}, \code{//docs:embedding},
2026.

\bibitem{runtime}
F.~Filmar and D.~Janković, ``The TxHDL runtime,'' project repository
\code{HDL/txhdl}, \code{//docs:runtime}, 2026.

\bibitem{axi}
F.~Filmar and D.~Janković, ``AXI in TxHDL,'' project repository
\code{HDL/txhdl}, \code{//docs:axi}, 2026.

\bibitem{examples}
F.~Filmar and D.~Janković, ``TxHDL by example,'' project repository
\code{HDL/txhdl}, \code{//docs:examples}, 2026.

\bibitem{cover}
F.~Filmar and D.~Janković, ``TxHDL documents,'' project repository
\code{HDL/txhdl}, \code{//docs:cover}, 2026.

\end{thebibliography}

\end{document}
